\documentclass[a4paper,10pt,final]{article}    %ngerman   % Neue Rechtschreibung, d.\,h. (Silbentrennung)
% -------------------------- Packages ---------------------------------------
\usepackage[utf8]{inputenc}        % UTF-8 input encoding (pdfLaTeX)
\usepackage[T1]{fontenc}           % Output font encoding
%\usepackage[a4paper,left=2.5cm,right=2.5cm,top=2.5cm,bottom=2.5cm]{geometry} % Margins
%\usepackage{helvet}                % Helvetica font
%\renewcommand{\familydefault}{\sfdefault} % Sans-serif default (Helvetica)
\usepackage{setspace}              % Line spacing
\onehalfspacing                    % 1.5-line spacing
\usepackage{graphicx}
\usepackage[skip=5pt]{caption}
\usepackage{etoolbox}
\AtBeginEnvironment{table}{\vspace{1em}}
\AtEndEnvironment{table}{\vspace{1em}}
\usepackage{float}
\usepackage{amsmath,amssymb}
\usepackage{booktabs}
\usepackage{pdflscape}
\usepackage{makecell}
\usepackage{subcaption}
\usepackage{natbib}                % American referencing style
\usepackage{hyperref}
\usepackage{authblk,xcolor}
\hypersetup{colorlinks=true,linkcolor=black,citecolor=black,urlcolor=blue}
\usepackage{threeparttable}   % for tablenotes
\usepackage{tabularx,ragged2e,makecell}
\renewcommand\tabularxcolumn[1]{>{\RaggedRight\arraybackslash}p{#1}}
\usepackage{siunitx}
\sisetup{
  round-mode     = places,
  round-precision= 3,
  table-number-alignment = center,
  detect-all
}

% --- Plots directory ---
\graphicspath{
  {../BFX_Results}
}

% -------------------------- Document ---------------------------------------

\setlength{\textwidth}{15.5cm} \setlength{\oddsidemargin}{.15cm} \addtolength{\parskip}{4pt}
\renewcommand{\baselinestretch}{1.4}
\textheight22cm
\topmargin-.25cm
\parindent.5cm

\begin{document}

\title{\large \bf Nonlinear Exchange Rate Dynamics and Behavioral Expectations: \\ New Evidence from Selected Emerging and Advanced Economies}

%\author{}
\author[a]{Lorenzo Biondi\thanks{Corresponding author. E-mail: {\tt lorenzobiondi.econ@gmail.com}. All errors remain ours. \newline The python code for replication is available at https://github.com/lorenzobiondiecon/str-buip-estimation}}
\author[b,c]{Christian R. Proa\~no}
\affil[a]{University of Siena, Italy}
\affil[b]{Otto-Friedrich-Universit\"at Bamberg, Bamberg, Germany}
\affil[c]{Centre for Applied Macroeconomic Analysis, Australian National University, Canberra, Australia}
\date{\today}
\maketitle

\begin{abstract}
    \noindent
    In this paper we estimate a behavioral uncovered interest rate parity (BUIP) in which exchange-rate returns are determined by the interest-rate differential (UIP) plus the market expectations on the future exchange rate, driven by two forecasting heuristics: chartist and fundamentalist based on a long-run purchasing power parity (PPP) fundamental. The endogenous market share of each rule follows a discrete-choice mechanism driven by relative performance of the rules, yielding nonlinear dynamics. Using monthly data on fourteen bilateral USD exchange rates over long samples ending in 2025, we compare a reduced-form logistic smooth-transition regression (STR) with the structural BUIP specification. The nonlinear models provide modest improvements in fit relative to a linear benchmark. Mixed evidence is found about the transition between regimes, with discontinuous threshold transitions more frequent than smooth functions. Fundamentalist forces are strongly identified by the BUIP model. Overall, in the data under study, allowing heterogeneous beliefs to enter market expectations partly helps to reconcile UIP with observed nonlinear exchange-rate dynamics while highlighting that PPP as the ``fundamental'' anchor remains an empirically relevant—but potentially market-dependent—choice.
\end{abstract}

\noindent \textbf{Keywords:} Exchange rates; Uncovered interest parity; Purchasing power parity; Heterogeneous agents; Behavioral expectations; Smooth transition regression; Nonlinear time series.

\noindent \textbf{JEL classifications:} F31; F32; G15; D84; C22; C52.

% 1. Journal of Evolutionary Economics

\thispagestyle{empty}
\clearpage

\tableofcontents

\thispagestyle{empty}
\clearpage

\pagenumbering{arabic}

%%%%%%%%%%%%%%%%%%%%%%%%%%%%%%%%%%%%%%%%%%%%%%%%%%%%%%%%%%%%%%%%%%%%%%%%%%%%%%%

\section{Introduction}
\label{chap:introduction}

Nearly half a century after the seminal contribution by \citet{MeeseRogoff1983}, nominal exchange rates remain notoriously hard to explain and forecast. Canonical parity conditions such as uncovered interest parity (UIP) and purchasing power parity (PPP) are routinely rejected, giving rise to the forward-premium anomaly and the PPP puzzle \citep{Rogoff1996,Engel2014}. A growing body of evidence attributes these anomalies to the interaction of heterogeneous market participants—most prominently, fundamentalists who anchor expectations to macroeconomic values and chartists who extrapolate recent trends \citep{FrankelFroot1987,BrockHommes1997,DeGrauwe2006}. Such behavioral models can reproduce excess volatility, fat-tailed returns, and regime-like swings, but robust empirical validation across currencies and long samples remains limited.

Understanding exchange-rate dynamics is central for monetary policy and macroprudential design. Persistent misalignments can erode competitiveness, while sudden depreciations can threaten financial stability—especially in economies with sizable foreign-currency liabilities. Linear parity regressions are ill suited to these features because adjustment pressures plausibly depend on the market state: arbitrage and mean reversion may be weak close to equilibrium but strengthen when deviations become large \citep{Sarno2006}. Nonlinear time-series frameworks, such as smooth-transition models, offer a reduced-form way to capture state dependence, while heterogeneous-agent models provide a structural mechanism—endogenous switching between forecasting heuristics—capable of generating nonlinear mean dynamics \citep{TaylorPeelSarno2001,DeGrauwe2006}. Empirically documenting how far such mechanisms can explain exchange-rate returns across a diverse set of currencies is therefore of direct interest to both international macroeconomics and policy.

Against this background, this paper tests UIP through the lens of a heterogeneous agent model with behavioral expectations developed by \citet{Proano2011,proano2013monetary}. We estimate a behavioral UIP (BUIP) law of motion in which monthly exchange-rate returns depend on the interest-rate differential and on market expectations formed by two forecasting rules: (i) a chartist trend-following rule and (ii) a fundamentalist rule that anchors expectations to a long run PPP fundamental. The market share of each rule evolves endogenously via a discrete-choice mechanism driven by relative performance, implying time variation in the strength of momentum versus PPP-consistent mean reversion \citep{BrockHommes1997,Proano2011,proano2013monetary}. To disentangle reduced-form state dependence from the structural switching implied by the theory, we estimate both (a) a logistic smooth-transition regression (STR), where regime changes are driven by observable transition variables (misalignment or volatility proxies), and (b) the structural BUIP specification, where the transition is endogenized through the profitability-based utilities of the two rules. Specifically, we investigate (i) whether the estimated transition function
captures the hypothesized switch between chartist and fundamentalist dominance; (ii) whether the
resulting regime-dependent coefficients possess the expected signs and statistical significance; and (iii)
whether systematic differences emerge between the economies under study.

We use monthly data for fourteen bilateral USD exchange rates—seven advanced economies (Australia, Canada, Euro Area, Japan, New Zealand, Switzerland, United Kingdom) and seven emerging economies (Brazil, Indonesia, Korea, Mexico, Philippines, Thailand, Türkiye)—with unbalanced samples extending back to the early 1980s and ending in 2025. Our results can be summarized as follows. First, the reduced-form STR indicates pervasive state dependence: linearity is rejected for 13 out of 14 currencies, and nonlinear specifications improve in-sample fit relative to a linear benchmark in terms of RMSE for 13 out of 14 countries. Nevertheless, once complexity is penalized, the evidence is weaker in several advanced economies, where the AIC frequently prefers the linear model, suggesting that nonlinearity is rewarded in emergent markets. Second, the STR delivers limited identification of the behavioral coefficients: statistically significant trend-following is found only for a small subset of currencies, and PPP-consistent mean reversion is detected only for a few cases, with some estimates implying opposite sign to the hypothesized one. Third, the structural BUIP specification yields a clearer decomposition into fundamentalist and chartist forces and identifies PPP-consistent mean reversion more frequently: the fundamentalist adjustment parameter is negative and statistically significant in eight currencies, while it is imprecisely estimated in several others and turns significantly positive in Mexico and T\"urkiye, suggesting that PPP is not an equally valid anchor across markets. Fourth, the chartist component is heterogeneous: some currencies display significant momentum, whereas others exhibit significant contrarian behavior at the monthly horizon. Overall, allowing heterogeneous expectations to enter UIP helps to rationalize nonlinearities in exchange-rate dynamics, but the empirical support for structural switching and for a unique fundamental anchor remains uneven across currency pairs.

The remainder of this paper is organized as follows. Section \ref{chap:literature} reviews theoretical and empirical work on exchange-rate puzzles, heterogeneous-agent models, and nonlinear econometric approaches. Section \ref{chap:methodology} presents the behavioral UIP framework and its econometric mapping and the estimation strategy. Section \ref{chap:results} presents the data and the empirical results. Section \ref{chap:conclusions} concludes and discusses limitations and possible extensions.

\section{Literature Review}
\label{chap:literature}

As is widely known, two benchmark parity conditions organize much of international finance: uncovered interest parity (UIP) and purchasing power parity (PPP), in absolute or relative form. Persistent empirical failures of these conditions---often summarized as ``puzzles'' (e.g., the exchange-rate disconnect, excess volatility, and the forward-premium/UIP anomaly)---have motivated a large literature that enriches the canonical macro view with nonlinear dynamics, heterogeneous trading behavior, and market microstructure mechanisms (see \citealp{SarnoTaylor2001} and \citealp{Evans2011}).

\paragraph{Classic monetary and portfolio-balance approaches.}
Early open-economy monetary models treat the exchange rate as an asset price pinned down by money-market equilibrium together with PPP and UIP. Three classic strands follow: flexible-price monetary models, sticky-price (overshooting) models, and portfolio-balance models. They differ mainly in the role of nominal rigidity and in whether domestic and foreign assets are perfect substitutes.
Flexible-price monetary models---e.g., \citep{Frenkel1976,Bilson1978}---offer a transparent link between exchange rates, monetary fundamentals, and expectations, recasting the exchange rate as the relative price of monies. They also highlight how credible monetary policy can stabilize the exchange rate. However, these models struggle to match high-frequency volatility and the weak short-run connection to fundamentals documented in the data.
Introducing nominal rigidities yields the canonical overshooting model of \citet{Dornbusch1976}. Because asset prices adjust quickly while goods prices adjust slowly, a permanent monetary expansion induces an immediate exchange-rate jump beyond its new long-run level; the exchange rate then gradually reverts as prices catch up. Overshooting is larger the greater nominal rigidity and the lower the interest sensitivity of money demand. \citet{Frankel1979} helps reconcile flexible- and sticky-price views by emphasizing real interest-rate differentials: short-run movements reflect liquidity and expectations effects, while long-run levels track monetary fundamentals. Empirically, \citet{Bilson1978} finds support for the monetary approach—using 1970s DM/\pounds{} data—underscoring the roles of expectations and speculation.
Relaxing perfect asset substitutability leads to portfolio-balance models, where domestic and foreign bonds are imperfect substitutes for risk-averse investors. A time-varying risk premium enters the parity condition and depends on relative asset supplies and wealth, feeding back into the exchange rate \citep{BransonHenderson1985,FrenkelMussa1985}. While attractive for policy analysis (e.g., sterilized interventions), subsequent evidence often points to signaling and coordination channels as more consequential than lasting portfolio effects \citep{SarnoTaylor2001}.

\paragraph{Micro-founded NOEM and present-value perspectives.}
In the mid-1990s, following the \citet{ObstfeldRogoff1995} Redux model, a new strand of micro-founded open-economy models with nominal rigidities emerged---the ``new open-economy macroeconomics'' (NOEM); see \citet{Lane2001}. Redux embeds overshooting in an intertemporal optimizing framework with current account dynamics and welfare analysis, yet key anomalies (e.g., the persistence of parity deviations) remain. Among later contributions, \citet{GaliMonacelli2005} develop a small-open-economy DSGE model with Taylor-rule monetary policy; under strong inflation targeting, the exchange rate can display near-random-walk behavior while still acting as a shock absorber.
A complementary present-value view links exchange rates to expected discounted future fundamentals: if $s_t$ equals the discounted value of future fundamentals plus a possible bubble component, then exchange rates may forecast future fundamentals rather than the reverse. Building on \citet{Lucas1982} and \citet{FrenkelMussa1985}, \citet{EngelWest2005} show that exchange rates forecast future macro fundamentals more effectively than fundamentals forecast exchange rates, consistent with high persistence in fundamentals translating into ``disconnect'' at short horizons.

\paragraph{Market microstructure and information aggregation.}
Macroeconomic models typically treat the exchange rate as an asset price determined by macro fundamentals. Market microstructure opens the trading ``black box'' of over-the-counter dealer markets where private information and inventory management matter. Order flow---the net of buyer- over seller-initiated trades---emerges as a proximate driver of dealers' quotes and transaction prices \citep{Evans2011}. Using high-frequency data, \citet{EvansLyons2002} show that order flow explains a substantial share of daily exchange-rate variation, consistent with dealers updating quotes in response to information conveyed by trades. This view explains seemingly ``newsless'' exchange-rate movements: the relevant information is in order flow rather than in public announcements \citep{Lyons2001}. Survey evidence aligns with this interpretation: \citet{CheungChinn2001} report that FX professionals view fundamentals as important at long horizons, while technical trading and bandwagon effects matter at short horizons.

\subsection{Heterogeneous-agent and behavioral exchange-rate models}
\label{subsec:ham}

\paragraph{Chartists vs.\ fundamentalists and endogenous switching.}
Building on standard DSGE/NOEM frameworks, a growing literature introduces heterogeneous expectations to capture the coexistence of distinct trading rules in FX markets. In the canonical setup, traders switch between a \emph{chartist} (trend-following, technical) rule and a \emph{fundamentalist} rule that anchors expectations to macroeconomic fundamentals. Following the discrete-choice switching mechanism of \citet{BrockHommes1997}, strategy weights evolve with relative past performance, yielding nonlinear aggregate dynamics with persistent swings and prolonged parity deviations, punctuated by abrupt corrections---an idea that goes back to \citet{FrankelFroot1990}.
Early contributions already emphasized that nonlinear interactions between chartists and fundamentalists can generate complex exchange-rate dynamics. \citet{de1993chaotic} embed these trading rules in a monetary model and show that the resulting nonlinear speculative dynamics can exhibit chaotic behavior, reproducing unit-root-like properties and biased forward premia even in the absence of explicit risk premia. Related work emphasizes ``animal spirits'' and sentiment dynamics: \citet{de2012animal} develop a model with evolutionary switching among optimistic and pessimistic beliefs (and, in extensions, fundamentalists and chartists), generating endogenous volatility, fat tails, and potentially chaotic dynamics. These models interpret excess volatility and ``disconnect'' as emergent outcomes of bounded rationality and strategy interaction rather than as failures of arbitrage per se.

\paragraph{Structural implementations and policy-relevant open-economy extensions.}
A first structural implementation on a DSGE skeleton is \citet{DeGrauwe2006}, who embed evolutionary switching in a dynamic macro model and reproduce momentum, excess volatility, bubbles, and crashes while preserving a long-run fundamentals anchor. In the same vein, \citet{Proano2011,proano2013monetary} augment a New-Keynesian DSGE core with a behavioral modification of UIP, combining IS, Phillips curve, and Taylor rule with an exchange-rate equation driven by heterogeneous expectations and switching.
Recent work extends heterogeneous-trader exchange-rate mechanisms into richer macro agent-based environments, explicitly linking FX dynamics to real activity, trade, and policy interventions. \citet{gori2018dynamic} propose a dynamic exchange-rate model with heterogeneous agents where the \emph{fundamental} exchange rate is endogenously determined by balance-of-payments conditions, thereby integrating heterogeneous speculation with real-side adjustment. Going further, \citet{delli2024complex} embed a currency market with chartists and fundamentalists into a two-country macro ABM with firm heterogeneity, innovation (R\&D), and international trade. In their framework, the exchange rate is jointly shaped by trade flows (real channel) and speculative demand for foreign currency (financial channel), and policy experiments with central-bank ``leaning against the wind'' highlight non-trivial trade-offs: aggressive interventions can compress speculative activity and cycle duration, but may also affect volatility in unintended ways.

\paragraph{Empirical evidence of heterogeneity and switching.}
A central challenge for heterogeneous-agent models is empirical validation. Time-series and survey evidence increasingly supports systematic heterogeneity and time variation in forecasting rules. \citet{manzan2007heterogeneous} estimate a structural chartist--fundamentalist model on monthly major USD exchange rates (1974--1998) and find strong in-sample fit for most currencies, with chartist behavior becoming stabilizing for sufficiently large deviations from fundamentals. Micro-level evidence using survey expectations is particularly informative. \citet{goldbaum2014empirical} use Consensus Economics forecasts to classify institutions by forecasting rule, document persistent heterogeneity across forecasters, and show that allowing for \emph{dynamic switching} materially improves fit relative to static classifications; moreover, switching is related to both period-specific and forecaster-specific factors.
Evidence also suggests that the intensity of switching may be weak or unstable in some contexts, raising identification and robustness issues. In a broader methodological contribution, \citet{kukacka2017estimation} propose nonparametric simulated maximum likelihood estimation for financial agent-based models and report, in an empirical application, that the switching coefficient can be statistically insignificant while trend-following behavior remains robust---a result consistent with the view that strategy heterogeneity is present even when evolutionary selection is difficult to pin down precisely in finite samples.

\paragraph{Estimation and identification in agent-based models.}
The recent literature has developed increasingly sophisticated estimation strategies for ABMs/HAMs, addressing concerns about overparameterization and nonlinear likelihoods. \citet{franke2011estimation} estimate a two-type (fundamentalist/chartist) agent-based asset-pricing model with \emph{structural stochastic volatility} using the method of simulated moments, obtaining a good fit for daily S\&P 500 returns and, in an FX application to USD--DEM, finding parameter changes consistent with a weaker role for chartism. Complementing simulation-based structural estimation, \citet{WesterhoffReitz2003} propose and estimate a nonlinear chartist--fundamentalist exchange-rate model using a STAR--GARCH specification on daily USD spot rates (against the DM, BP, and JPY), providing empirical support for regime-dependent dynamics consistent with heterogeneous trading and nonlinear mean reversion, while also questioning the presumption that fundamentalists are always stabilizing. Together with \citet{kukacka2017estimation}, this line of work helps move heterogeneous-agent models from qualitative ``stylized-fact'' narratives toward structurally estimated and statistically testable frameworks.

\subsection{Nonlinear dynamics and empirical evidence beyond linear parity models}
\label{subsec:nonlinear}

\paragraph{Nonlinearity from limits to arbitrage and state dependence.}
A key implication of heterogeneous-agent and behavioral models is \emph{state-dependent} dynamics. Near equilibrium, chartist and fundamentalist pressures may coexist and small shocks can be amplified if technical trading becomes temporarily more profitable; as misalignment grows, mean-reversion strategies become relatively more profitable, inducing a regime shift and accelerating correction. This mechanism yields nonlinear adjustment consistent with long stretches of near-random-walk behavior interspersed with sharp corrections.
Empirical support for nonlinear exchange-rate dynamics also comes from smooth transition and regime-switching time-series models. Foundational treatments of STR/STAR modeling include \citet{GrangerTerasvirta1993,FransesVanDijk2000,Terasvirta1994,vanDijk2002,GrangerTerasvirta2010,TsayChen2019}. For PPP, a central result is nonlinear mean reversion in real exchange rates: adjustment is weak near equilibrium but accelerates when misalignment is large. In an ESTAR model, \citet{Michael1997} attribute this to transaction costs creating a band of inaction; \citet{TaylorPeelSarno2001} similarly document nonlinear mean reversion and reject the random walk for real exchange rates.
Markov-switching models provide an alternative reduced-form representation of persistent exchange-rate ``swings.'' \citet{EngelHamilton1990} show that regime-dependent specifications can outperform the random-walk benchmark for the U.S.\ dollar, consistent with latent-state shifts in conditional drifts.

\paragraph{Chaos and heterogeneous beliefs.}
Beyond generic nonlinear adjustment, several contributions explicitly investigate chaotic dynamics implied by heterogeneous beliefs. Building on the theoretical possibility established by \citet{de1993chaotic}, \citet{federici2012euro} estimate a continuous-time heterogeneous-beliefs model for the EUR/USD (with fundamentalists, chartists, and other actors) and conduct a battery of chaos diagnostics (phase plots, correlation dimension, Lyapunov exponents). While the empirical relevance of strict chaos remains debated, these studies reinforce the broader point that economically motivated heterogeneity can induce intrinsic nonlinearities that are observationally close to stochastic processes.
Recent work also shows that simple agent-based implementations of chartist--fundamentalist interaction can match multiple FX stylized facts in modern samples. \citet{mignot2025explaining} propose an agent-based version of de Grauwe’s chaotic exchange-rate model and demonstrate that it reproduces bubbles/crashes, excess volatility, fat-tailed returns, and volatility clustering for major currency pairs using daily data (1999--2022). Such evidence supports the view that heterogeneous trading rules can generate empirically realistic dynamics even with parsimonious behavioral primitives.

\subsection{Fundamentals, misalignment, and the horizon-dependent role of macro variables}
\label{subsec:fundamentals}

A recurring empirical regularity is that fundamentals matter more at long horizons than at short horizons. Microstructure evidence interprets the short-run disconnect partly as an informational issue: dispersed beliefs and private signals are aggregated through trading (order flow), while public macro announcements explain only a small share of daily movements \citep{Lyons2001,EvansLyons2002,EvansLyons2008}. Survey evidence is consistent with horizon dependence: \citet{CheungChinn2001} report that professionals emphasize fundamentals at long horizons and technical/bandwagon effects at short horizons.
Heterogeneous-agent models provide a structural rationale for this pattern. If chartists dominate near equilibrium or in tranquil periods, the exchange rate can drift away from fundamentals with limited immediate correction; when misalignment becomes sufficiently large, fundamentalists become more influential and mean-reversion strengthens. Empirically, \citet{manzan2007heterogeneous} and \citet{goldbaum2014empirical} support heterogeneity and switching in forecasting behavior, while macro ABM extensions emphasize that fundamentals may be endogenously linked to trade flows and real activity \citep{gori2018dynamic,delli2024complex}. Overall, the literature suggests that fundamentals remain an essential anchor, but their effect is nonlinear and horizon-dependent rather than constant and contemporaneous.

\subsection{Uncovered interest parity and the evolving evidence on the ``Fama puzzle''}
\label{subsec:uip}

UIP asserts that the expected depreciation equals the interest rate differential (under rational expectations and a zero risk premium), implying a unit slope in the Fama regression. Empirically, short-horizon tests typically reject UIP and often produce negative slope estimates (the forward-premium anomaly). More recent evidence stresses that the relationship is \emph{time-varying} and sensitive to monetary policy regimes and the level of interest rates.
\citet{bussiere2018new} document that the Fama coefficient changed sign around the global financial crisis and then flipped again, and argue that these reversals are not well explained by standard risk proxies such as the VIX. Importantly, when replacing ex post depreciation with survey-based expectations, UIP regressions exhibit fewer apparent breaks and expected depreciation is much more consistently positively correlated with interest differentials, pointing to a central role for expectation errors and learning rather than a stable risk-premium story alone.
Focusing on the U.S.\ dollar, \citet{engel2022reconsideration} show that the predictability of excess returns from interest differentials becomes weak when interest rates are near the zero lower bound and that the relationship is unstable over time. They also report that inflation differentials provide more stable predictive content for future excess returns than interest differentials, consistent with delayed market adjustment to predictable monetary-policy changes.
At the same time, there is renewed evidence that UIP may hold \emph{in the long run} even if it fails at short horizons. \citet{baillie2025yes} test long-horizon UIP using Durbin regression methods designed to improve efficiency in the presence of serial correlation and dynamics, and conclude that while UIP is rejected at short horizons, the long-run restriction implied by UIP (and the long-run International Fisher relation) is not rejected for major currencies over extended samples. Taken together, the recent literature suggests a nuanced view: UIP is a poor short-run description of exchange-rate dynamics, but its long-run content may be closer to the data once time variation, policy regimes, and expectation formation are properly accounted for---a perspective that aligns naturally with heterogeneous-expectations models in which carry-trade profitability and deviations from UIP are state-dependent.

The literature shows a clear structural narrative. Different expectations and shifting conditions can create nonlinear exchange-rate dynamics. They can also explain the varying roles of fundamentals and reshape the empirical content of UIP.
However, identifying these effects is tough. Reduced-form nonlinear models may show state dependence but don’t clarify the economic mechanisms. Structural switching models can struggle if parameters are weakly identified or if the chosen fundamental proxy is incorrect. This need drives our empirical strategy.
Next, we derive the behavioral UIP law of motion. We map it into (i) a reduced-form LSTR specification and (ii) a structural BUIP model with profitability-based switching. We then estimate both frameworks across fourteen USD currency pairs using long monthly samples that end in 2025.

\section{Methodology}
\label{chap:methodology}

\subsection{The Underlying Theoretical Framework}
\label{chap:theory}

The underlying theoretical framework for our empirical analysis builds on \citet{Proano2011,proano2013monetary}, who considered the exchange rate determination in a small open economy with perfect capital mobility and boundedly rational expectations.

Let $s_t$ denote the log nominal exchange rate (domestic price of foreign currency); $p_t$ and $p_t^{*}$ the log domestic and foreign price levels; and $\pi_t \equiv p_t - p_{t-1}$ and $\pi_t^{*} \equiv p_t^{*} - p_{t-1}^{*}$ the respective inflation rates. Short-term nominal interest rates in the domestic and foreign economies are $i_t$ and $i_t^{*}$, respectively.

FX traders are assumed to choose between two forecasting heuristics for the future nominal exchange rate: A fundamentalist rule and a chartist (trend-following) rule. Following \citet{Proano2011,proano2013monetary}, their decisions -- as well as those of the central bank and the real sector -- are based on the information set available up to period $t-1$.

Fundamentalists expect mean reversion of the nominal exchange rate toward its long-run value $f$, which is assumed to be consistent with the PPP in absolute form, i.e.
\[
    f_t \equiv p_t - p_t^{*}.
\]
Accordingly, the fundamentalists' expected nominal exchange is given by
\begin{eqnarray}
    E_t^{f}[s_{t+1}]
    &=& s_{t-1} + \beta_s^{f}\bigl(f_{t-1}-s_{t-1}\bigr) \nonumber\\
    &=& s_{t-1} + \beta_s^{f}\bigl(p_{t-1} - p_{t-1}^{*} -s_{t-1}\bigr),\nonumber\\
    &=& s_{t-1} - \beta_s^{f}\,\eta_{t-1},
    \label{eq:f_expect}
\end{eqnarray}
where $\beta_s^{f}$ is a positive parameter measuring the speed of adjustment of the exchange rate to the fundamental value assumed by the fundamentalist traders, and
\[
    \eta_t \equiv s_t + p_t^{*} - p_t = s_t - f_t,
\]
is the log real exchange rate, which is also a measure of misalignment from the PPP long run value of the nominal rate.

By contrast, chartists extrapolate last period's change to form their expectation:
\begin{equation}
    E_t^{c}[s_{t+1}]
    = s_{t-1} + \beta_s^{c}\,\Delta s_{t-1},
    \label{eq:c_expect}
\end{equation}
where $\Delta s_{t-1} \equiv s_{t-1}-s_{t-2}$ and $\beta_s^{c}$ is a positive parameter measuring the intensity of trend-following behavior of chartists.

Following the discrete-choice framework of \citet{BrockHommes1997,BrockHommes1998}, the market share (``market mood'') of fundamentalists, $\omega_t \in (0,1)$, depends on the last-period relative profitability of the two rules. Let $\psi_{t-1}^{j}$ denote the period $t-1$ earnings from rule $j\in\{f,c\}$ and $\sigma^{2}_{j,t-1}$ its squared forecast error. In \citet{Proano2011,proano2013monetary} the weight is
\begin{equation}
    \omega_t
    = \frac{\exp\!\bigl[\gamma\,(\psi^{f}_{t-1}-\sigma^{2}_{f,t-1})\bigr]}
    {\exp\!\bigl[\gamma\,(\psi^{f}_{t-1}-\sigma^{2}_{f,t-1})\bigr]
        + \exp\!\bigl[\gamma\,(\psi^{c}_{t-1}-\sigma^{2}_{c,t-1})\bigr]},
    \qquad \gamma>0,
    \label{eq:omega}
\end{equation}
so that $\omega_t \to 1$ when the fundamentalist rule dominates and $\omega_t \to 0$ when technical trading dominates.

Last-period earnings of investing one unit of domestic currency in the foreign currency, $\psi^{j}_{t-1}$, depend on the accuracy implied by the rule and the realized excess return:\footnote{Here we deviate from the original formulation of \citet{Proano2011,Proano2011} and replace the realized profit component ($S_{t-1}(1+i^) - S_{t-2}(1+i)$) used there with the foreign-exchange excess return -- also interpretable as the UIP deviation or FX risk premium --- following \citet{Engel2014}.}
\begin{equation}
    \psi^{j}_{t-1}
    = \bigl[i^*_{t-1} - i_{t-1} + \Delta s_{t-1} \bigr]
    \operatorname{sgn}\bigl[E^{j}_{t-2}\Delta s_{t-1}\bigr]
    \label{eq:psi}
\end{equation}
where $S_t \equiv e^{s_t}$ is the nominal exchange rate in levels, and $\operatorname{sgn}(E^{j}_{t-2}\Delta s_t)\in\{-1,0,1\}$ represents the anticipated appreciation/depreciation suggested by rule $j$ at $t-2$. Each rule captures that profits are positive (negative) when the rule correctly (incorrectly) predicts appreciation/depreciation. Traders choose by comparing the rules' utilities associated with each trading rule defined as
\[
    U_{f,t}=\psi_{f,t} - \sigma^2_{f,t},
\]
\[
    U_{c,t} = \psi_{c,t} - \sigma^2_{c,t},
\]
with $\sigma_j$ being the squared error of the logarithmic forecast\footnote{While \citet{Proano2011,proano2013monetary} used the squared forecast error of the exchange rate level $\bigl(E^{j}_{t-2}S_{t-1}-S_{t-1}\bigr)^{2}$, this formulation facilitates identification by consistently working in logarithmic terms, thereby avoiding scale inconsistencies across variables, which is an important consideration given the sensitivity of the logistic function to the magnitude of its argument.} given by rule $j$:
\begin{equation}
    \sigma^{2}_{j,t-1} = (E_{t-2}^j s_{t-1} - s_{t-1})^2,
    \label{eq:sigma}
\end{equation}
This formulation mirrors the behavioral heterogeneous-agents literature \citep{DeGrauwe2006,Hommes2013}.

Further, since fundamentalist strategies require access to information and the associated costs of processing it, they are assumed to be more costly than chartist strategies. To capture this, we introduce a cost parameter $k > 0$ that enters negatively in the fundamentalist utility:
\[
    U_{f,t}= \psi_{f,t} - \sigma^2_{f,t} -k,
\]
so that the market share function takes the form
\begin{equation}
    \label{eq:omega_k}
    \omega_t = \frac{1}{1+exp[-\gamma ((\psi_{f,t} - \sigma^2_{f,t}) - (\psi_{c,t} - \sigma^2_{c,t})  -k)]}
\end{equation}

This formulation resembles a smooth transition regression (STR) with an estimable threshold. The threshold $k$ facilitates identification and has a clear economic interpretation: the utility differential must exceed a certain level before fundamentalist traders become the majority and potentially dominate market dynamics. Since such episodes occur less frequently, the stabilizing, mean-reverting influence of fundamentalist strategies is correspondingly reduced, thereby amplifying destabilizing dynamics. Among the specifications considered, this adjustment proved to be the most effective in terms of identification, yielding empirically reasonable results.

With perfect capital mobility, uncovered interest parity (UIP) holds \emph{ex ante} at the market level, but with behavioral expectations. Define the market expectation as the weight-averaged forecast,
\begin{equation}
    E_t^{m}[s_{t+1}] \equiv \omega_t E_t^{f}[s_{t+1}] + (1-\omega_t) E_t^{c}[s_{t+1}].
    \label{eq:Em}
\end{equation}
Thus, the exchange rate updates based on the expectations of traders made at the beginning of t for the following period and on the contemporaneous interest rate differential. The equation thus, takes the form of a Behavioral UIP \citep{proano2013monetary}:
\begin{equation}
    s_t = i_t^{*}-i_t + E_t^{m}[s_{t+1}],
    \label{eq:uip_level}
\end{equation}
which, after subtracting $s_{t-1}$ and substituting \eqref{eq:f_expect}--\eqref{eq:c_expect}, yields the behaviorally founded law of motion of the exchange rate:
\begin{equation}
    \Delta s_t
    = (i_t^{*}-i_t)
    - \omega_t\,\beta_s^{f}\,\eta_{t-1}
    + (1-\omega_t)\,\beta_s^{c}\,\Delta s_{t-1}.
    \label{eq:beh-uip}
\end{equation}
Equation \eqref{eq:beh-uip} shows that exchange-rate returns depend on the interest-rate differential and on two state-dependent behavioral forces: a PPP-correcting term whose strength grows with $\omega_t$ and a momentum term whose strength grows with $1-\omega_t$. Because $\omega_t$ is a nonlinear function of lagged profitability \eqref{eq:omega}, \eqref{eq:beh-uip} is a nonlinear difference equation capable of generating persistent swings and regime-like behavior consistent with empirical evidence on nonlinear mean reversion and deviations from parity \citep{TaylorPeelSarno2001,DeGrauwe2006,MenkhoffTaylor2007}.

\label{sec:real}

To connect FX dynamics to macroeconomic stabilization, \citet{Proano2011,proano2013monetary} embeds \eqref{eq:beh-uip} in a parsimonious real-side block featuring an output gap equation, a Phillips-Curve relationship and a different Taylor-like interest rate rules. A higher trend-following intensity $\beta_s^{c}$ raises persistence in $\Delta s_t$ and can amplify the real effects of shocks indirectly via $\eta_{t-1}$ in the IS equation of \citet{proano2013monetary}. The sensitivity parameter $\gamma$ controls how sharply the market shifts between regimes; larger $\gamma$ increases the amplitude of fluctuations in $\omega_t$ but does not by itself imply macro instability if policy satisfies the Taylor principle and real-side parameters remain in the stable region \citep{proano2013monetary}. These mechanisms echo the broader heterogeneous-agents literature, in which rule switching generates bubbles, crashes, and prolonged deviations from parity while maintaining a fundamentals-based anchor in the long run \citep{DeGrauwe2006,ManzanWesterhoff2007}. Empirically, nonlinear mean reversion in real exchange rates \citep{TaylorPeelSarno2001,MichaelNobayPeel1997} and the widespread use of technical analysis \citep{MenkhoffTaylor2007,ParkIrwin2007,CheungChinn2001} motivate the coexistence of rules and the profitability-based switching assumed here. The model is also consistent with the observed weak link between macro news and high-frequency FX moves \citep{EvansLyons2008,EvansLyons2002,Evans2011}.


\subsection{Econometric Methodology}
\label{chap:econometrics}

Equation~\eqref{eq:beh-uip} fits in the Smooth Transition Regression class of econometric models, first developed by \citet{GrangerTerasvirta1993,Terasvirta1994}, and popularized in modelling nonlinear time series and exchange rate dynamics with the works of \citet{TaylorPeelSarno2001,vanDijk2002,GrangerTerasvirta2010}. These statistical models allow two or more states of the world and enable dynamic switching among them through a transition function. Most of the works in the literature rely on the Logistic (LSTR) or Exponential (ESTR) functions to define the switching behavior. The present work first estimates the LSTR model as defined below to understand the relevance of such a modeling strategy. Then the full behavioral law of motion with the utilities of the two types of traders is estimated via NLS to understand (i) whether model adds something to the ``plain'' STR specification, and (ii) whether the bilateral exchange rates dynamics in question can be resolved in part through this modelling choice.
The reduced form LSTR model is also useful to understand which variables are relevant to define the switching mechanism, providing useful insights about market expectations.

Our aim is to capture the departures of the exchange rate returns from the UIP, so we define the dependent variable $y$ as the difference between the monthly log return and the interest rate differential, i.e.,
$y = \Delta s_{t} - (i^* - i)$. This formulation preserves the UIP restriction while attributing short-horizon departures to expectation formation---an interpretation aligned with the microstructure framework \citep{Evans2011} and consistent with recent evidence supporting long-horizon UIP validity \citep{baillie2025yes}. Two regressors capture, respectively, trend-following and fundamental mean reversion trading strategies: the lagged return $\Delta s_{t-1}$ and the lagged real exchange rate $\eta_{t-1}=s_{t-1}+p^*_{t-1}-p_{t-1}$. Coefficients evolve smoothly between two regimes through the logistic function
\[
    G\!\left(z_{t-d};\gamma ,c\right)
    =\Bigl[1+\exp\!\bigl\{-\gamma\,(z_{t-d}-c)\bigr\}\Bigr]^{-1},
\]
where $\gamma>0$ measures the switching intensity and $c$ is the threshold value of the transition variable $z_{c,t-d}$, the driver of the regime-changes. To allow for regime switching, the candidate transition variables are measures of exchange-rate misalignment from fundamentals, all expressed in absolute value so that the switching mechanism depends on magnitude rather than sign and avoids one-sided thresholds, which would have no clear economic interpretation\footnote{Another possibility would have been estimating a LSTR2 model, with two threshold values $c_1,c_2$, such that for values inside the band $[c_1,c_2]$ the chartist regime dominates, while for values outside the band the fundamentalist regime takes over. We chose the simpler LSTR specification for parsimony and comparison with the structural BUIP model.}.

The candidates are---all in absolute value: (i)$|\eta_t|$, the real exchange rate ; (ii) $|\Delta s_t|$, the nominal return; (iii) $|i_t^{*}-i_t|$, the interest-rate differential; (iv) $|f_t|$ the PPP fundamental value; (v) $|\Delta f_t|$ the fundamental value PPP change; (vi) $|\Delta f_t - \Delta s_t|$, the real exchange rate returns; and (vii) $|\Delta s_t - i_t^{*}-i_t|$ the excess returns. Each series is considered up to lag $t-5$, and the transition variable is selected via an $LM$ test for linearity as described below.

The estimated parameter $c$ represents the threshold value beyond which fundamentalist traders dominate market dynamics. Nevertheless, given the smooth nature of the transition function, both types of traders are almost always present, with pure dominance occurring only in the limiting cases where the transition function $G(\cdot)$ approaches 0 or 1. The threshold thus denotes the point beyond which the proportion of fundamentalist traders exceeds 50\% of the market.

Stacking these elements yields, for every economy,
\begin{equation}
    \label{eq:str_econometric}
    y_t
    = \alpha_0 + \beta^c \Delta s_{t-1} + G_t \bigl(\alpha_1 + \beta^f \eta_{t-1} - \beta^c \Delta s_{t-1}\bigr)
    +\varepsilon_t,
\end{equation}
with $\alpha_0$ and $\alpha_1$ the intercepts in the two regimes and $\varepsilon_{t}\sim\text{i.i.d.}(0,\sigma^{2})$. This formulation facilitates a comprehensive understanding of the distinctions between the two regimes. Within the chartist regime, where $G = 0$, the model is defined as:

$$y_t = \alpha_0 + \beta^c \Delta s_{t-1}$$

Conversely, in the fundamentalist regime, where $G = 1$, the model transitions to:

$$y_t = \alpha_0 + \alpha_1 + \beta^f \eta_{t-1}$$. A negative  $\beta^f$ should indicate mean reversion towards the PPP fundamental value, while a positive $\beta^c$ should capture momentum. The chartist coefficient could capture the destabilizing effect of technical trading if the estimated value is greater than 1.

A more structural approach estimates \eqref{eq:beh-uip} while \emph{building} $\omega_t$ from the discrete-choice mapping \eqref{eq:omega}. The estimable system is
\begin{equation}
    y_t
    = \alpha_{0,UIP} + \beta_{c,UIP} \Delta s_{t-1} + \omega_t (\alpha_{1,UIP} + \beta_{f,UIP} \eta_{t-1} - \beta_{c,UIP} \Delta s_{t-1}),\label{eq:buip_econometric}
\end{equation}
with $\omega_t$ defined as in \eqref{eq:omega_k}:
\begin{equation}
    \label{eq:fund_market_share}
    \omega_t = \frac{1}{1+exp[-\gamma (U_{f,t} - U_{c,t})]}.
\end{equation}

The structural approach tightens the link with theory and estimates directly endogenous regime switching, while LSTR offers a flexible reduced form that can capture state dependence and interpreted through the lens of \citet{proano2013monetary} model. The empirical chapter will implement both, beginning with LSTR for parsimony and moving to structural NLS for completeness and robustness.

The econometric estimation routine follows the methodology of \citet{GrangerTerasvirta2010}, encompassing model specification, estimation, and evaluation.\paragraph{Specification}The specification phase entails testing the null hypothesis of linearity against the Smooth Transition Regression ($STR$) alternative. This procedure is implemented via a Lagrange Multiplier ($LM$) test derived from a Taylor expansion of the transition function around the parameter $\gamma = 0$ \citep{Terasvirta1994, vanDijk2002}. The test is conducted for each candidate transition variable and associated lag (up to a five-month horizon). The transition variable yielding the lowest $p$-value, which indicates a rejection of the null hypothesis at the highest significance level, is selected for the estimation stage. The standard procedure for model type selection (LSTR vs. LSTR2 vs. ESTR) is ignored in this analysis, as the functional form is dictated by the underlying theoretical model.

\paragraph{Estimation}

Equations~\ref{eq:str_econometric} and \ref{eq:buip_econometric} are fitted by NonLinear Least Squares (NLS). The procedure is slightly different for the two models.
\begin{itemize}
    \item \textbf{STR}: Because the likelihood surface can feature several local maxima, \citet{GrangerTerasvirta2010} suggest beginning with a grid search over different combinations of $(\gamma,c)$ for every candidate transition variable. At each grid point the linear coefficients conditional on $(\gamma,c)$ are obtained by Ordinary Least Squares (OLS) and the best estimation is selected by minimizing the residual sum of squares (RSS). The resulting $(\hat{\gamma}, \hat{c})$ are used as starting values for the NLS regression, performed via the Trust-Region optimization algorithm (a variant of the Levenberg-Marquardt algorithm that allows to set bounds), that updates all parameters jointly, subject to the positivity constraint on $\gamma$. Convergence tolerances on parameters and objective are set to $10^{-8}$, and the iteration cap is fixed at 20000. Identification of $\gamma$ is notoriously difficult as noted by \citep{GrangerTerasvirta2010}, but a grid search over $\gamma$ can help reach a local maximum. Adhering to the methodology proposed by \citet{GrangerTerasvirta2010}, $\gamma$ is normalized by the standard deviation of the transition variable. This procedure reduces scale sensitivity and enhances the identification properties of the parameter.
    \item \textbf{BUIP}: The estimation strategy parallels that of the STR model; however, the conditional ordinary least squares step cannot be implemented due to the functional specification of the market share $\omega_t$. Consequently, we initiate the procedure by fitting a linear benchmark specification ($y_t = const + \beta_c \Delta s_{t-1} + \beta_f \eta_{t-1}$) through ordinary least squares. The estimated coefficients serve to construct approximate utilities $U_{f,t}$ and $U_{c,t}$, under the assumption that the cost parameter equals zero ($k=0$). A grid of candidate starting values for the parameter pair $(\gamma,k)$ is constructed based on quantiles of the utility differential $(U_{f,t}-U_{c,t})$. For each grid point, the corresponding market share $\omega_t$ is computed. Given these market share proxies, the model becomes linear in the remaining parameters conditional on $(\gamma,k)$, enabling estimation via ordinary least squares. These preliminary parameter estimates subsequently serve as initial values for the final nonlinear least squares optimization.
\end{itemize}

Reported in the tables are the confidence intervals and significance levels based on heteroskedasticity and autocorrelation consistent (HAC) standard errors computed via the Newey-West estimator with a lag length selected according to the \citet{newey1994automatic} rule.

\paragraph{Evaluation}
The first step in the evaluation phase involves checking if the estimated threshold parameter ($\hat{c}$ in the STR and $\hat{k}$ in the BUIP model) lies withing the range of the transition variable and is statistically significant. This provides a first indication of a correct identification of non linear models. The countries for which the estimated threshold lies outside the range are treated as non-identified for the sample-data under study, and excluded from further analysis.
Then, residual diagnostics include Ljung-Box tests for autocorrelation, ARCH Lagrange-multiplier tests for conditional heteroskedasticity, Jarque-Bera tests for normality.

The next section presents the results of the empirical analysis.

\section{Results}
\label{chap:results}

\subsection{Data Description}
\label{chap:data}

We employ data obtained from the IMF International Financial Statistics (IFS) and OECD Main Economic Indicators. The dataset comprises monthly observations for fourteen economies: Australia (AU), Brazil (BRA), Canada (CAN), the Euro Area (EA), Indonesia (ID), Japan (JPN), Korea (KOR), Mexico (MEX), New Zealand (NZ), the Philippines (PH), Switzerland (CH), Thailand (TH), Türkiye (TR), and the United Kingdom (GBR). We exploit the maximum available sample period for each economy, resulting in an unbalanced panel spanning from as early as the 1980s to the latest data available in 2025. Appendix~\ref{app:data_details} provides a detailed overview of the sample periods for each country.

\begin{table}[htbp]
    \centering
    \caption{First Four Moments of Monthly Log Returns}
    \label{tab:desc_stats}
    \begin{tabular}{lrrrrrr}
        \toprule
                       & $Mean_{\Delta s}$ & $Var_{\Delta s}$ & $Skew_{\Delta s}$ & $Kurt_{\Delta s}$ & $Mean_{f}$ & $Var_{f}$ \\
        Country        &                   &                  &                   &                   &            &           \\
        \midrule
        Australia      & 0.0008            & 0.0010           & 0.7045            & 2.5387            & 0.3314     & 0.1274    \\
        Brazil         & 0.0488            & 0.0104           & 1.6606            & 3.1150            & 3.3373     & 0.8037    \\
        Canada         & 0.0002            & 0.0004           & 0.4249            & 3.9327            & -0.0051    & 0.1346    \\
        Euro Area      & -0.0000           & 0.0007           & 0.1595            & 1.2867            & -0.0919    & 0.2374    \\
        Indonesia      & 0.0057            & 0.0037           & 5.3678            & 69.4273           & 1.3997     & 0.1515    \\
        Japan          & -0.0010           & 0.0009           & -0.2226           & 1.3990            & -2.0937    & 0.4672    \\
        Korea          & 0.0012            & 0.0009           & 0.7188            & 12.5920           & 0.5421     & 0.3243    \\
        Mexico         & 0.0116            & 0.0024           & 6.0905            & 68.4200           & 1.8292     & 0.6185    \\
        New Zealand    & -0.0005           & 0.0012           & 0.3714            & 1.2110            & -0.1009    & 0.0610    \\
        Philippines    & 0.0038            & 0.0008           & 3.5908            & 25.9394           & 1.0881     & 0.2147    \\
        Switzerland    & -0.0019           & 0.0010           & -0.1025           & 0.9630            & -0.3797    & 0.4980    \\
        Thailand       & 0.0007            & 0.0008           & 0.2737            & 27.8384           & 0.2695     & 0.1638    \\
        Türkiye        & 0.0234            & 0.0025           & 2.4488            & 15.0845           & 2.5129     & 0.2761    \\
        United Kingdom & 0.0005            & 0.0008           & 0.1674            & 2.4225            & 0.0776     & 0.0722    \\
        \bottomrule
    \end{tabular}
\end{table}

Three series are required per country: the nominal spot rate $S_t$ expressed as the domestic price of US dollars, the consumer-price index $P_t$, and the short term interest rate $i_t$ (which is almost equal to the monetary policy rate). US variables complete the construction of each country's dataset for real exchange rates and interest differentials construction. All variables -- except for the interest rate -- are used in natural logarithms ($s_t = \ln S_t$, $p_t = \ln P_t^{\text{DOM}}$, $p_t^{*} = \ln P_t^{\text{US}}$). Monthly log returns are computed as  the differences of the log nominal exchange rate ($\Delta s_t = s_t - s_{t-1}$). Interest rates are converted from annual percentage rates to effective monthly decimals via $(1+\text{APR})^{1/12}-1$.

\begin{figure}[htbp]
    \centering
    \includegraphics[width=\textwidth]{log_nominal_vs_ppp_6panel.pdf}
    \caption{Nominal exchange rate $s_t$ (solid line) and long run PPP fundamental $f_t$ (dashed line) for selected economies.}
    \label{fig:log_nominal_vs_ppp_6panel}
\end{figure}
Table~\ref{tab:desc_stats} shows marked heterogeneity across currencies: emerging markets (e.g., Brazil, Mexico, Indonesia) exhibit higher means, variances, and extreme skewness/kurtosis, while advanced economies display near-zero means, low dispersion, and thinner tails; the PPP fundamental $f_t$ are modestly positive in mean for most emerging markets---consistent with the depreciation trend of emerging-markets currencies---and near zero or negative for advanced peers.

Figure~\ref{fig:log_nominal_vs_ppp_6panel} displays the log nominal exchange rate $s_t$ (solid black line) and the PPP long run fundamental value $f_t$ (dashed gray line) for selected economies over the sample period. The variables have been normalized to the first observation, thus assuming that PPP holds at that date. The vertical distance between the nominal exchange rate and the PPP fundamental measures the misalignment from purchasing power parity. It's immediately clear the presence of the Disconnect Puzzle \citep{ObstfeldRogoff2000} in each economy, where each series show large deviations from the absolute PPP fundamental, that persist for long periods---sometimes more than 15 years---with short run movements always present. The difference between emergent and advanced economies is striking. The emergent countries show a historical trend of depreciation against the US dollar, while the advanced countries fluctuate around an almost stable mean. The deviations from the PPP long run value are present in all countries.

Plots of the exchange-rate returns $\Delta s_t$ and the change in the PPP fundamental $\Delta f_t$ (not shown) confirm the higher volatility of the nominal exchange rate compared to the PPP changes \citep{EngelWest2005}. The emerging markets show higher volatility levels, as confirmed by the results in Table~\ref{tab:desc_stats}.


\subsection{Estimation Results}
\label{sec:baseline_all_z}

\subsubsection{Empirical Results of the STR Estimation}
\label{sec:str_results}

The STR estimations provide evidence that exchange-rate dynamics can display state-dependent behavior, consistent with the nonlinear specification introduced in Section~\ref{chap:methodology}. As a preliminary step, we test the null of linearity against the STR alternative and select, for each country, the transition variable that yields the strongest evidence of nonlinearity. Table~\ref{tab:transition_variables} reports the selected transition variable and the corresponding LM linearity-test $p$-value. Linearity is rejected at conventional levels for 13 out of 14 economies. Australia is the only case in which the linearity null cannot be rejected(p-value=0.2426), suggesting that an STR-type regime mechanism is not empirically supported for the AUS/USD currency. This is confirmed by the identification results reported in \ref{tab:validation_summary} in Appendix~\ref{app:validation_summary}, where the estimated threshold lies outside the range of the transition variable, indicating a failure to identify the model. Along with Australia, also Mexico, Thailand, Turkyie and the UK show $\hat c$ estimates outside the range of the transition variable. While linearity is still rejected for these countries, this results suggests that the identification problem resides in the estimation algorithm. While Australia results won't be discussed further, these other countries' results are reported for completeness, but will be treated with caution.

Beyond the widespread rejection of linearity, the selected transition variables indicate heterogeneity in the signals that govern regime changes. For a large subset of countries (Brazil, Canada, Indonesia, Japan, Korea, Mexico), the preferred transition variable is a measure of (absolute) relative misalignment, consistent with switching driven by the magnitude of deviations from fundamentals. In other cases (Switzerland, Thailand, Turkiye, UK) the transition variable is the fundamental value change in absolute value. This cross-country heterogeneity is consistent with the view that nonlinear switching is a robust feature, driven mostly by fundamental misalignment, but that the precise mechanism varies across currency markets.

\begin{table}[htbp]
    \caption{Selected Transition Variables for STR Models (LM Linearity Test)}
    \label{tab:transition_variables}
    \centering{
    \begin{tabular}{lll}
        \toprule
        Country        & Transition Variable                 & P-value \\
        \midrule
        Australia      & $|\Delta f_{t-3} - \Delta s_{t-3}|$ & 0.2426  \\
        Brazil         & $|\Delta f_{t-2} - \Delta s_{t-2}|$ & 0.0000  \\
        Canada         & $|\Delta f_{t-2} - \Delta s_{t-2}|$ & 0.0000  \\
        Euro Area      & $|\Delta s_{t-3}|$                  & 0.0112  \\
        Indonesia      & $|\Delta f_{t-4} - \Delta s_{t-4}|$ & 0.0000  \\
        Japan          & $|\Delta f_{t-2} - \Delta s_{t-2}|$ & 0.0071  \\
        Korea          & $|\Delta f_{t-4} - \Delta s_{t-4}|$ & 0.0000  \\
        Mexico         & $|\Delta f_{t-4} - \Delta s_{t-4}|$ & 0.0000  \\
        New Zealand    & $|(i^{*}_{t-4} - i_{t-4})|$         & 0.0059  \\
        Philippines    & $|\Delta s_{t-3}|$                  & 0.0173  \\
        Switzerland    & $|\Delta f_{t-2}|$                  & 0.0000  \\
        Thailand       & $|\Delta f_{t-5}|$                  & 0.0000  \\
        Türkiye        & $|\Delta f_{t-4}|$                  & 0.0000  \\
        United Kingdom & $|\Delta f_{t-3}|$                  & 0.0151  \\
        \bottomrule
    \end{tabular}}
    \begin{tablenotes}[para,flushleft]
        \small
        \textit{Note}: Transition-variable definitions follow the notation introduced in Section~\ref{chap:methodology}.
    \end{tablenotes}
\end{table}

Table~\ref{tab:str_results_ci} reports the STR coefficient estimates (with 95\% confidence intervals). Positive and significant $\beta_c$ estimates are obtained for Brazil, Canada, and Japan only, indicating that the reduced form model fails to capture simple trend-following behavior. New Zealand instead features a significantly negative $\beta_c$, consistent with contrarian technical trading within the sample under analysis. These results suggest that either the model specification is too simple to capture technical trading dynamics in FX markets, or that technical trading is not a dominant force in the countries under consideration.

Turning to the fundamentalist component, $\beta_f$, the estimates are more selective. A negative $\beta_f$ implies mean reversion towards the fundamental PPP value. Statistically significant mean reversion is found for Brazil, Korea, New Zealand, and Switzerland, while the Euro Area displays weaker evidence (significant at the 10\% level). By contrast, Mexico exhibits a positive and significant $\beta_f$, which is inconsistent with fundamental reversion hypothesized and points to an identification issue---confirmed by the negative value of the estimated $c$, which is inconsistent with the specification of the transition variables, which are all $\geq 0$. Poor identification of the fundamentalist coefficient suggests that this reduced-form STR model may not fully capture the underlying behavioral dynamics. This motivates research in the direction of either different reduced-form specifications---e.g. different fundamentals, more complex structures---or estimating the full structural behavioral UIP model, as done in Section~\ref{sec:buip_results}.

\begin{table}[htbp]
    \centering
    \caption{STR Model Estimation Results with 95\% Confidence Intervals}
    \label{tab:str_results_ci}
    \resizebox{\textwidth}{!}{%
        \begin{tabular}{l|cccccc}
            \toprule
            Country        & $\alpha_0$                    & $\alpha_1$ & $\beta_c$ & $\beta_f$ & $\gamma$ & $c$ \\
            \midrule
            Australia      & \makecell{15.1129                                                                   \\ \scriptsize{(-746274.8410, 746305.0668)}} & \makecell{-15.1073 \\ \scriptsize{(-751014.4889, 750984.2743)}} & \makecell{294.2944 \\ \scriptsize{(-6424922.9100, 6425511.4987)}} & \makecell{-0.0179 \\ \scriptsize{(-0.1297, 0.0938)}} & \makecell{0.4071 \\ \scriptsize{(-1.6624, 2.4766)}} & \makecell{-0.4191 \\ \scriptsize{(-1835.0218, 1834.1837)}} \\
            Brazil         & \makecell{0.0139                                                                    \\ \scriptsize{(-0.0147, 0.0425)}} & \makecell{\textbf{0.3476}*** \\ \scriptsize{(0.1059, 0.5893)}} & \makecell{\textbf{1.9925}*** \\ \scriptsize{(1.3599, 2.6251)}} & \makecell{\textbf{-0.2821}** \\ \scriptsize{(-0.5168, -0.0474)}} & \makecell{\textbf{0.9070}** \\ \scriptsize{(0.1987, 1.6154)}} & \makecell{\textbf{0.1037}*** \\ \scriptsize{(0.0385, 0.1689)}} \\
            Canada         & \makecell{-0.0002                                                                   \\ \scriptsize{(-0.0026, 0.0022)}} & \makecell{0.0056 \\ \scriptsize{(-0.0011, 0.0123)}} & \makecell{\textbf{0.1816}** \\ \scriptsize{(0.0273, 0.3359)}} & \makecell{-0.0187* \\ \scriptsize{(-0.0402, 0.0028)}} & \makecell{256.0000 \\ \scriptsize{(-8867.5062, 9379.5062)}} & \makecell{\textbf{0.0113}*** \\ \scriptsize{(0.0102, 0.0124)}} \\
            Euro Area      & \makecell{-0.0059                                                                   \\ \scriptsize{(-0.0135, 0.0016)}} & \makecell{0.0029 \\ \scriptsize{(-0.0058, 0.0116)}} & \makecell{0.0840 \\ \scriptsize{(-0.1256, 0.2937)}} & \makecell{-0.0224* \\ \scriptsize{(-0.0469, 0.0021)}} & \makecell{256.0000 \\ \scriptsize{(-7218.4533, 7730.4533)}} & \makecell{\textbf{0.0036}*** \\ \scriptsize{(0.0023, 0.0049)}} \\
            Indonesia      & \makecell{\textbf{0.0091}***                                                        \\ \scriptsize{(0.0064, 0.0118)}} & \makecell{0.7043 \\ \scriptsize{(-0.5728, 1.9814)}} & \makecell{0.4063 \\ \scriptsize{(-0.1821, 0.9948)}} & \makecell{-0.0733 \\ \scriptsize{(-0.2074, 0.0608)}} & \makecell{198.1245 \\ \scriptsize{(-8527.5461, 8923.7950)}} & \makecell{\textbf{0.0722}*** \\ \scriptsize{(0.0578, 0.0865)}} \\
            Japan          & \makecell{-0.0020                                                                   \\ \scriptsize{(-0.0052, 0.0012)}} & \makecell{-0.0408 \\ \scriptsize{(-0.1111, 0.0295)}} & \makecell{\textbf{0.2274}*** \\ \scriptsize{(0.1115, 0.3433)}} & \makecell{0.0086 \\ \scriptsize{(-0.0068, 0.0239)}} & \makecell{512.0001 \\ \scriptsize{(-1179116.8173, 1180140.8175)}} & \makecell{0.0279 \\ \scriptsize{(-0.2503, 0.3061)}} \\
            Korea          & \makecell{\textbf{0.0035}**                                                         \\ \scriptsize{(0.0007, 0.0063)}} & \makecell{\textbf{3.0003}*** \\ \scriptsize{(1.2617, 4.7389)}} & \makecell{0.2684 \\ \scriptsize{(-0.0952, 0.6320)}} & \makecell{\textbf{-0.4215}*** \\ \scriptsize{(-0.6598, -0.1831)}} & \makecell{1.2871 \\ \scriptsize{(-0.5169, 3.0912)}} & \makecell{\textbf{0.0822}** \\ \scriptsize{(0.0129, 0.1514)}} \\
            Mexico         & \makecell{72.8714                                                                   \\ \scriptsize{(-189412.4816, 189558.2243)}} & \makecell{-73.0846 \\ \scriptsize{(-193414.4166, 193268.2475)}} & \makecell{751.3892 \\ \scriptsize{(-1671013.0987, 1672515.8771)}} & \makecell{\textbf{0.0814}** \\ \scriptsize{(0.0187, 0.1440)}} & \makecell{0.9064 \\ \scriptsize{(-0.3385, 2.1512)}} & \makecell{-0.3324 \\ \scriptsize{(-74.5624, 73.8975)}} \\
            New Zealand    & \makecell{0.0005                                                                    \\ \scriptsize{(-0.0029, 0.0040)}} & \makecell{\textbf{0.0351}*** \\ \scriptsize{(0.0189, 0.0513)}} & \makecell{\textbf{-0.1446}** \\ \scriptsize{(-0.2557, -0.0336)}} & \makecell{\textbf{-0.0659}*** \\ \scriptsize{(-0.0984, -0.0334)}} & \makecell{256.0000 \\ \scriptsize{(-1403.0448, 1915.0448)}} & \makecell{\textbf{0.0023}*** \\ \scriptsize{(0.0023, 0.0024)}} \\
            Philippines    & \makecell{\textbf{0.0295}***                                                        \\ \scriptsize{(0.0084, 0.0506)}} & \makecell{0.0023 \\ \scriptsize{(-0.0471, 0.0517)}} & \makecell{-0.0288 \\ \scriptsize{(-0.1729, 0.1154)}} & \makecell{-0.0066 \\ \scriptsize{(-0.0185, 0.0052)}} & \makecell{5544.9990 \\ \scriptsize{(-112876.6736, 123966.6717)}} & \makecell{\textbf{0.0004}*** \\ \scriptsize{(0.0003, 0.0004)}} \\
            Switzerland    & \makecell{\textbf{-0.0038}***                                                       \\ \scriptsize{(-0.0065, -0.0012)}} & \makecell{\textbf{0.0348}*** \\ \scriptsize{(0.0119, 0.0576)}} & \makecell{0.0769* \\ \scriptsize{(-0.0057, 0.1594)}} & \makecell{\textbf{-0.2274}*** \\ \scriptsize{(-0.3809, -0.0740)}} & \makecell{63.8486 \\ \scriptsize{(-1925.3333, 2053.0306)}} & \makecell{\textbf{0.0078}*** \\ \scriptsize{(0.0077, 0.0079)}} \\
            Thailand       & \makecell{-10.6922                                                                  \\ \scriptsize{(-261308.0737, 261286.6893)}} & \makecell{23.9689 \\ \scriptsize{(-484762.3210, 484810.2588)}} & \makecell{0.3122 \\ \scriptsize{(-1169.1484, 1169.7728)}} & \makecell{-0.0660 \\ \scriptsize{(-303.7298, 303.5977)}} & \makecell{0.0002 \\ \scriptsize{(-4.7013, 4.7018)}} & \makecell{5.0255 \\ \scriptsize{(-131877.1827, 131887.2337)}} \\
            Türkiye        & \makecell{-153.8334                                                                 \\ \scriptsize{(-246735.8948, 246428.2281)}} & \makecell{153.9550 \\ \scriptsize{(-241156.2837, 241464.1936)}} & \makecell{N.I.} & \makecell{-0.0009 \\ \scriptsize{(-0.0207, 0.0190)}} & \makecell{0.3841 \\ \scriptsize{(-0.0750, 0.8433)}} & \makecell{N.I.} \\
            United Kingdom & \makecell{0.0092                                                                    \\ \scriptsize{(-0.0136, 0.0320)}} & \makecell{-436.1348 \\ \scriptsize{(-16680975.5908, 16680103.3212)}} & \makecell{0.0808 \\ \scriptsize{(-0.0767, 0.2382)}} & \makecell{-527.8929 \\ \scriptsize{(-20190372.3510, 20189316.5653)}} & \makecell{0.2400 \\ \scriptsize{(-0.5780, 1.0579)}} & \makecell{0.1272 \\ \scriptsize{(-465.5153, 465.7697)}} \\
            \bottomrule
        \end{tabular}%
    }
    \begin{tablenotes}[para,flushleft]
        \small
        \textit{Note}: 95\% confidence intervals shown in parentheses below parameter estimates.
        Significance levels: *** $p<0.01$, ** $p<0.05$, * $p<0.10$.
        Bold estimates are significant at 5\% level.
        $\alpha_0$, $\alpha_1$ = intercepts; $\beta_c$ = chartist coefficient;
        $\beta_f$ = fundamentalist coefficient; $\gamma$ = transition speed; $c$ = threshold.
    \end{tablenotes}
\end{table}

The switching parameter provides additional insights but also highlight identification constraints that are common in STR applications. For the countries in which the smoothness parameter is precisely estimated, the implied transition is economically plausible. Brazil and Korea display moderate $\gamma$ estimates, consistent with gradual regime switching, as confirmed in Figure~\ref{fig:str_scatter_4x4}. For several other countries, $\gamma$ is estimated at (or near) large values, indicating a threshold switching rather than a smooth one\footnote{As stressed in \citet{GrangerTerasvirta2010}, the t-value of the gamma estimate does not have the usual interpretation. Large values suggest a TR (Threshold Regression) model where the transition function is discontinuous, and the model jumps from one regime to the other rather than switching gradually as with the STR.}.

The threshold parameter $c$ is precisely estimated in many cases (e.g., Brazil, Canada, Euro Area, Indonesia, New Zealand, Philippines, Switzerland, Korea). Brazil, Indonesia and Korea estimates indicate that fundamental expectations start dominating the market when the real exchange rate returns---which measures the misalignment of the returns and the PPP change---get to 7-10\%. However, Indonesia's $\beta_f$ is imprecisely estimated. This is consistent with the narrative that in emergent economies fundamentalists traders take more time to intervene and markets are more unstable, due to destabilizing forces of strong technical trading.

Residual diagnostics for the STR models are reported in Table~\ref{tab:str_diagnostics}. As typical for high-frequency FX returns, the Jarque--Bera test strongly rejects normality across all countries. The Ljung--Box results suggest that for several advanced economies (e.g., Australia, Euro Area, Switzerland, United Kingdom) serial correlation is not present, while most emerging markets still exhibits residual autocorrelation (notably Brazil, Indonesia, Korea, Mexico, Philippines, Thailand). ARCH effects remain widespread, indicating that the STR specification captures conditional mean nonlinearities but does not model conditional heteroskedasticity; an STR--GARCH extension is therefore a natural next step.

\begin{table}[htbp]
    \caption{STR Model Diagnostic Tests (Test Statistic and p-value)}
    \label{tab:str_diagnostics}
    \begin{tabular}{llll}
        \toprule
                       & Ljung-Box Q                & ARCH LM                    & Jarque-Bera                  \\
        Country        &                            &                            &                              \\
        \midrule
        Australia      & 8.2801 (0.6015)            & \textbf{14.9009} (0.0049)  & \textbf{196.5628} (0.0000)   \\
        Brazil         & \textbf{59.2758} (0.0000)  & \textbf{39.6404} (0.0000)  & \textbf{2545.6774} (0.0000)  \\
        Canada         & \textbf{23.4180} (0.0093)  & \textbf{15.2635} (0.0042)  & \textbf{276.5862} (0.0000)   \\
        Euro Area      & 8.3538 (0.5943)            & \textbf{48.8331} (0.0000)  & \textbf{18.8333} (0.0001)    \\
        Indonesia      & \textbf{44.2145} (0.0000)  & \textbf{50.3805} (0.0000)  & \textbf{73729.2951} (0.0000) \\
        Japan          & \textbf{20.1975} (0.0274)  & 5.3142 (0.2566)            & \textbf{28.9678} (0.0000)    \\
        Korea          & \textbf{38.7903} (0.0000)  & \textbf{162.5248} (0.0000) & \textbf{1867.5550} (0.0000)  \\
        Mexico         & \textbf{184.5590} (0.0000) & 7.4513 (0.1139)            & \textbf{11919.5325} (0.0000) \\
        New Zealand    & \textbf{19.1951} (0.0379)  & \textbf{31.3239} (0.0000)  & \textbf{73.7937} (0.0000)    \\
        Philippines    & \textbf{25.1323} (0.0051)  & 8.9700 (0.0619)            & \textbf{9863.3155} (0.0000)  \\
        Switzerland    & 10.3420 (0.4110)           & 6.0242 (0.1973)            & \textbf{18.0328} (0.0001)    \\
        Thailand       & \textbf{19.3667} (0.0358)  & \textbf{52.8779} (0.0000)  & \textbf{15587.9390} (0.0000) \\
        Türkiye        & 14.5205 (0.1505)           & \textbf{14.0027} (0.0073)  & \textbf{8272.7336} (0.0000)  \\
        United Kingdom & 4.6851 (0.9112)            & \textbf{15.0545} (0.0046)  & \textbf{90.7274} (0.0000)    \\
        \bottomrule
    \end{tabular}
\end{table}


\begin{figure}[htbp]
    \centering

    % ===== Row 1 =====
    \begin{subfigure}{0.24\textwidth}
        \centering
        \includegraphics[width=\linewidth]{Australia_str_scatter.pdf}
    \end{subfigure}\hfill
    \begin{subfigure}{0.24\textwidth}
        \centering
        \includegraphics[width=\linewidth]{Brazil_str_scatter.pdf}
    \end{subfigure}\hfill
    \begin{subfigure}{0.24\textwidth}
        \centering
        \includegraphics[width=\linewidth]{Canada_str_scatter.pdf}
    \end{subfigure}\hfill
    \begin{subfigure}{0.24\textwidth}
        \centering
        \includegraphics[width=\linewidth]{Euro_Area_str_scatter.pdf}
    \end{subfigure}

    \medskip

    % ===== Row 2 =====
    \begin{subfigure}{0.24\textwidth}
        \centering
        \includegraphics[width=\linewidth]{Indonesia_str_scatter.pdf}
    \end{subfigure}\hfill
    \begin{subfigure}{0.24\textwidth}
        \centering
        \includegraphics[width=\linewidth]{Japan_str_scatter.pdf}
    \end{subfigure}\hfill
    \begin{subfigure}{0.24\textwidth}
        \centering
        \includegraphics[width=\linewidth]{Korea_str_scatter.pdf}
    \end{subfigure}\hfill
    \begin{subfigure}{0.24\textwidth}
        \centering
        \includegraphics[width=\linewidth]{Mexico_str_scatter.pdf}
    \end{subfigure}

    \medskip

    % ===== Row 3 =====
    \begin{subfigure}{0.24\textwidth}
        \centering
        \includegraphics[width=\linewidth]{New_Zealand_str_scatter.pdf}
    \end{subfigure}\hfill
    \begin{subfigure}{0.24\textwidth}
        \centering
        \includegraphics[width=\linewidth]{Philippines_str_scatter.pdf}
    \end{subfigure}\hfill
    \begin{subfigure}{0.24\textwidth}
        \centering
        \includegraphics[width=\linewidth]{Switzerland_str_scatter.pdf}
    \end{subfigure}\hfill
    \begin{subfigure}{0.24\textwidth}
        \centering
        \includegraphics[width=\linewidth]{Thailand_str_scatter.pdf}
    \end{subfigure}

    \medskip

    % ===== Row 4 =====
    \begin{subfigure}{0.24\textwidth}
        \centering
        \includegraphics[width=\linewidth]{Türkiye_str_scatter.pdf}
    \end{subfigure}\hfill
    \begin{subfigure}{0.24\textwidth}
        \centering
        \includegraphics[width=\linewidth]{United_Kingdom_str_scatter.pdf}
    \end{subfigure}\hfill
    \begin{subfigure}{0.24\textwidth}
        \centering
        % Empty
    \end{subfigure}\hfill
    \begin{subfigure}{0.24\textwidth}
        \centering
        % Empty
    \end{subfigure}
    \caption{STR model: scatter plots of transition function $G_t$ against transition variable $z_t$.}
    \label{fig:str_scatter_4x4}
\end{figure}


\subsubsection{Estimation of the Behavioral UIP Model}
\label{sec:buip_results}

Table~\ref{tab:buip_results_ci} reports the estimates of the structural BUIP model of \citet{Proano2011,proano2013monetary} defined in Equations~\eqref{eq:buip_econometric} and \eqref{eq:fund_market_share}. Relative to the reduced-form STR, this specification imposes an explicit evolutionary switching mechanism through the utility differential. The structural model also improves interpretation of the results: coefficients can be read as the strength of chartist versus fundamentalist forces, while the switching parameters describe how strongly relative profitability translates into market-share dynamics (or so-called ``animal spirits'' in the spirit of \citet{de2012animal}).


We first note that the estimated $k$ cost parameter is identified outside of the interval of the utility differential for Japan, Korea and the United Kingdom, suggesting that ill identification for these countries. The rest of the estimated $k$ indicate a positive significance cost for Brazil, Canada, Euro Area and T\"urkyie, while a negative significance cost for Indonesia, New Zealand, Philippines, showing heterogeneity in the endogenous selecting mechanism of trading strategies across countries.

The fundamentalist parameter $\beta_f$ is negative and statistically significant in a broad subset of countries (Australia, Euro Area, Indonesia, New Zealand, Philippines, Switzerland, Thailand), supporting the presence of stabilizing mean-reversion forces to the fundamental value when fundamental trading becomes relatively more profitable than chartist trading. At the same time, $\beta_f$ is not universally identified: it is imprecise for Brazil, Canada, Japan, and the United Kingdom. Moreover, Mexico and Türkiye display a positive and significant $\beta_f$, implying that, under the hypothesized definition of the fundamental forecasting strategy, the implied ``fundamentalist'' term does not correspond to stabilizing reversion in those markets. This pattern reinforces the interpretation that the empirical relevance of a given fundamental anchor can be market-dependent and time-varying, and that structural identification remains sensitive to how the fundamental forecast is specified.

\begin{table}[htbp]
    \centering
    \caption{BUIP Model Estimation Results with 95\% Confidence Intervals}
    \label{tab:buip_results_ci}
    \resizebox{\textwidth}{!}{%
        \begin{tabular}{l|cccccc}
            \toprule
            Country        & $\alpha_0$                   & $\alpha_1$      & $\beta_c$        & $\beta_f$ & $\gamma$ & $k$ \\
            \midrule
            Australia      & \makecell{0.0455                                                                               \\ \scriptsize{(-0.3391, 0.4302)}} & \makecell{-0.0382 \\ \scriptsize{(-0.4220, 0.3457)}} & \makecell{0.0123 \\ \scriptsize{(-2.7092, 2.7339)}} & \makecell{\textbf{-0.0159}** \\ \scriptsize{(-0.0316, -0.0001)}} & \makecell{812.7878 \\ \scriptsize{(-2535968.5719, 2537594.1475)}} & \makecell{-0.1377 \\ \scriptsize{(-5.8102, 5.5347)}} \\
            Brazil         & \makecell{\textbf{0.0460}***                                                                   \\ \scriptsize{(0.0301, 0.0618)}} & \makecell{-0.0698 \\ \scriptsize{(-0.1987, 0.0592)}} & \makecell{\textbf{1.7350}*** \\ \scriptsize{(1.5062, 1.9638)}} & \makecell{0.0063 \\ \scriptsize{(-0.0312, 0.0439)}} & \makecell{\textbf{2.3011}** \\ \scriptsize{(0.3223, 4.2798)}} & \makecell{\textbf{0.0998}** \\ \scriptsize{(0.0152, 0.1844)}} \\
            Canada         & \makecell{0.0019*                                                                              \\ \scriptsize{(-0.0002, 0.0040)}} & \makecell{-0.0284 \\ \scriptsize{(-0.0735, 0.0168)}} & \makecell{\textbf{-0.1349}** \\ \scriptsize{(-0.2392, -0.0307)}} & \makecell{-0.0297 \\ \scriptsize{(-0.1852, 0.1259)}} & \makecell{2.3282 \\ \scriptsize{(-0.4906, 5.1469)}} & \makecell{\textbf{0.0725}*** \\ \scriptsize{(0.0400, 0.1049)}} \\
            Euro Area      & \makecell{0.0012                                                                               \\ \scriptsize{(-0.0022, 0.0045)}} & \makecell{\textbf{-0.0266}*** \\ \scriptsize{(-0.0360, -0.0172)}} & \makecell{\textbf{0.0053}** \\ \scriptsize{(0.0006, 0.0100)}} & \makecell{\textbf{-0.0935}*** \\ \scriptsize{(-0.1427, -0.0444)}} & \makecell{206.7932 \\ \scriptsize{(-5561.4961, 5975.0825)}} & \makecell{\textbf{0.0278}*** \\ \scriptsize{(0.0155, 0.0401)}} \\
            Indonesia      & \makecell{\textbf{1.8482}***                                                                   \\ \scriptsize{(1.6995, 1.9968)}} & \makecell{\textbf{-1.4450}*** \\ \scriptsize{(-1.6514, -1.2386)}} & \makecell{\textbf{-4.3036}*** \\ \scriptsize{(-4.8034, -3.8038)}} & \makecell{\textbf{-0.0429}*** \\ \scriptsize{(-0.0674, -0.0184)}} & \makecell{\textbf{2.2241}** \\ \scriptsize{(0.4633, 3.9848)}} & \makecell{\textbf{-0.5137}*** \\ \scriptsize{(-0.8090, -0.2184)}} \\
            Japan          & \makecell{1.0760                                                                               \\ \scriptsize{(-6.9135, 9.0655)}} & \makecell{-1.0686 \\ \scriptsize{(-9.0535, 6.9162)}} & \makecell{56.7780 \\ \scriptsize{(-346.2992, 459.8552)}} & \makecell{-0.0027 \\ \scriptsize{(-0.0136, 0.0082)}} & \makecell{0.0164 \\ \scriptsize{(-0.2089, 0.2417)}} & \makecell{-14.2687 \\ \scriptsize{(-228.0732, 199.5359)}} \\
            Korea          & \makecell{N.I.}              & \makecell{N.I.} & \makecell{1.9776                              \\ \scriptsize{(-30287.2782, 30291.2334)}} & \makecell{\textbf{-0.0352}*** \\ \scriptsize{(-0.0544, -0.0161)}} & \makecell{7.9629 \\ \scriptsize{(-4176869.9632, 4176885.8891)}} & \makecell{N.I.} \\
            Mexico         & \makecell{\textbf{0.0234}***                                                                   \\ \scriptsize{(0.0179, 0.0290)}} & \makecell{\textbf{-0.0761}*** \\ \scriptsize{(-0.1279, -0.0243)}} & \makecell{\textbf{0.7313}*** \\ \scriptsize{(0.4473, 1.0153)}} & \makecell{\textbf{0.0259}*** \\ \scriptsize{(0.0070, 0.0449)}} & \makecell{299.4982 \\ \scriptsize{(-8724.3951, 9323.3916)}} & \makecell{0.0041 \\ \scriptsize{(-0.0019, 0.0101)}} \\
            New Zealand    & \makecell{\textbf{0.0100}***                                                                   \\ \scriptsize{(0.0026, 0.0175)}} & \makecell{0.0100 \\ \scriptsize{(-0.0026, 0.0225)}} & \makecell{\textbf{-0.2628}*** \\ \scriptsize{(-0.3855, -0.1402)}} & \makecell{\textbf{-0.0393}*** \\ \scriptsize{(-0.0611, -0.0174)}} & \makecell{619.0809 \\ \scriptsize{(-5548.9511, 6787.1129)}} & \makecell{\textbf{-0.0014}*** \\ \scriptsize{(-0.0023, -0.0004)}} \\
            Philippines    & \makecell{\textbf{0.0081}***                                                                   \\ \scriptsize{(0.0050, 0.0112)}} & \makecell{\textbf{0.6574}*** \\ \scriptsize{(0.6124, 0.7023)}} & \makecell{\textbf{0.0308}*** \\ \scriptsize{(0.0180, 0.0436)}} & \makecell{\textbf{-0.1746}*** \\ \scriptsize{(-0.1894, -0.1599)}} & \makecell{33.2207 \\ \scriptsize{(-96.7996, 163.2410)}} & \makecell{\textbf{-0.3544}*** \\ \scriptsize{(-0.4267, -0.2821)}} \\
            Switzerland    & \makecell{-0.0022                                                                              \\ \scriptsize{(-0.0062, 0.0018)}} & \makecell{0.0075 \\ \scriptsize{(-0.0021, 0.0171)}} & \makecell{0.0240 \\ \scriptsize{(-0.0411, 0.0891)}} & \makecell{\textbf{-0.0817}*** \\ \scriptsize{(-0.1255, -0.0379)}} & \makecell{9.7171 \\ \scriptsize{(-14.1794, 33.6136)}} & \makecell{0.0021 \\ \scriptsize{(-0.0090, 0.0131)}} \\
            Thailand       & \makecell{0.1490                                                                               \\ \scriptsize{(-0.4401, 0.7380)}} & \makecell{-0.0328 \\ \scriptsize{(-0.6886, 0.6230)}} & \makecell{-0.0044 \\ \scriptsize{(-0.0324, 0.0235)}} & \makecell{\textbf{-0.0343}*** \\ \scriptsize{(-0.0547, -0.0140)}} & \makecell{0.7657 \\ \scriptsize{(-1.8699, 3.4013)}} & \makecell{-0.2108 \\ \scriptsize{(-0.8474, 0.4257)}} \\
            Türkiye        & \makecell{\textbf{0.0400}***                                                                   \\ \scriptsize{(0.0332, 0.0468)}} & \makecell{\textbf{-0.0505}*** \\ \scriptsize{(-0.0708, -0.0302)}} & \makecell{\textbf{0.5654}*** \\ \scriptsize{(0.3731, 0.7578)}} & \makecell{\textbf{0.0361}*** \\ \scriptsize{(0.0177, 0.0544)}} & \makecell{1553.3028 \\ \scriptsize{(-3954.4186, 7061.0242)}} & \makecell{\textbf{0.0016}*** \\ \scriptsize{(0.0009, 0.0023)}} \\
            United Kingdom & \makecell{0.0009                                                                               \\ \scriptsize{(-0.0015, 0.0032)}} & \makecell{2.5795 \\ \scriptsize{(-241.7991, 246.9581)}} & \makecell{\textbf{0.0006}*** \\ \scriptsize{(0.0002, 0.0010)}} & \makecell{-0.0595 \\ \scriptsize{(-0.4613, 0.3422)}} & \makecell{3.8513 \\ \scriptsize{(-2.1357, 9.8383)}} & \makecell{0.1443 \\ \scriptsize{(-0.8335, 1.1220)}} \\
            \bottomrule
        \end{tabular}%
    }
    \begin{tablenotes}[para,flushleft]
        \small
        \textit{Note}: 95\% confidence intervals shown in parentheses below parameter estimates.
        Significance levels: *** $p<0.01$, ** $p<0.05$, * $p<0.10$.
        Bold estimates are significant at 5\% level.
        $\alpha_0$, $\alpha_1$ = intercepts; $\beta_c$ = chartist coefficient;
        $\beta_f$ = fundamentalist coefficient; $\gamma$ = transition speed; $c$ = cost/threshold parameter.
    \end{tablenotes}
\end{table}
The chartist coefficient $\beta_c$ is precisely estimated in many cases and displays substantial cross-country variation in sign. Positive and significant trend-following (chartist) force is found for Brazil, Euro Area, Mexico, Philippines and T\"urkyie, while negative and significant contrarian trading forces characterizes Canada, Indonesia, and New Zealand. This heterogeneity suggests that rather than simple trend following, technical trading may encompass a variety of strategies. This suggests once again that time variation and market structure may play a role in shaping the trading strategies of the market participants.

The switching intensity $\gamma$ is statistically significant and economically plausible for Brazil and Indonesia. For other countries, $\gamma$ is often weakly identified\footnote{See footnote 4}. Figure~\ref{fig:buip_scatter_4x4} presents the scatter plots of the transition function against the utility differential. The plots show a smoother transition function for Brazil, Canada and Switzerland, while a threshold like function for the rest of the countries.
The cost/threshold parameter $k$ is generally within the empirical support of the utility differential, but there are notable exceptions (Japan, Korea, United Kingdom) where $k$ lies outside the observed range; these cases should be interpreted with caution and might imply non-identification of the non-linear model as pointed out earlier. This result is confirmed by Figure~\ref{fig:buip_scatter_4x4} where the three countries display an invariant transition function.

Diagnostic tests for the BUIP residuals (Table~\ref{tab:buip_diagnostics}) mirror the STR diagnostics: non-normality is pervasive and ARCH effects remain common. The Ljung--Box test suggests that the BUIP specification reduces serial correlation for some developed markets (e.g., Australia, Canada, Euro Area, Switzerland, United Kingdom), while several emerging markets still display residual autocorrelation, indicating that additional dynamics (e.g., regime changes, risk premia, or time-varying volatility) may be needed for a fuller description.

\begin{table}
    \caption{BUIP Model Diagnostic Tests (Test Statistic and p-value)}
    \label{tab:buip_diagnostics}
    {\centering
    \begin{tabular}{llll}
        \toprule
                       & Ljung-Box Q                & ARCH LM                    & Jarque-Bera                  \\
        Country        &                            &                            &                              \\
        \midrule
        Australia      & 6.9792 (0.7274)            & \textbf{16.1039} (0.0029)  & \textbf{95.8223} (0.0000)    \\
        Brazil         & \textbf{116.5459} (0.0000) & \textbf{42.9387} (0.0000)  & \textbf{1137.8206} (0.0000)  \\
        Canada         & 13.6909 (0.1876)           & \textbf{15.7658} (0.0034)  & \textbf{459.3778} (0.0000)   \\
        Euro Area      & 13.1940 (0.2130)           & \textbf{44.5248} (0.0000)  & \textbf{34.9023} (0.0000)    \\
        Indonesia      & \textbf{18.4401} (0.0480)  & \textbf{12.1655} (0.0162)  & \textbf{14776.0672} (0.0000) \\
        Japan          & \textbf{18.9679} (0.0407)  & 5.8822 (0.2081)            & \textbf{34.0992} (0.0000)    \\
        Korea          & \textbf{23.6200} (0.0087)  & \textbf{118.8522} (0.0000) & \textbf{2303.7557} (0.0000)  \\
        Mexico         & \textbf{73.4924} (0.0000)  & 7.3842 (0.1169)            & \textbf{76830.4361} (0.0000) \\
        New Zealand    & \textbf{19.0202} (0.0400)  & \textbf{35.3253} (0.0000)  & \textbf{57.6146} (0.0000)    \\
        Philippines    & \textbf{44.5250} (0.0000)  & 8.7590 (0.0674)            & \textbf{17646.0861} (0.0000) \\
        Switzerland    & 9.9739 (0.4428)            & \textbf{17.5657} (0.0015)  & \textbf{12.9100} (0.0016)    \\
        Thailand       & 17.6522 (0.0611)           & \textbf{58.8718} (0.0000)  & \textbf{12733.6133} (0.0000) \\
        Türkiye        & \textbf{42.9295} (0.0000)  & \textbf{43.7378} (0.0000)  & \textbf{4739.0173} (0.0000)  \\
        United Kingdom & 5.2295 (0.8753)            & 8.9690 (0.0619)            & \textbf{101.2747} (0.0000)   \\
        \bottomrule
    \end{tabular}}
\end{table}

\begin{figure}[htbp]
    \centering

    % ===== Row 1 =====
    \begin{subfigure}{0.24\textwidth}
        \centering
        \includegraphics[width=\linewidth]{Australia_buip_scatter.pdf}
    \end{subfigure}\hfill
    \begin{subfigure}{0.24\textwidth}
        \centering
        \includegraphics[width=\linewidth]{Brazil_buip_scatter.pdf}
    \end{subfigure}\hfill
    \begin{subfigure}{0.24\textwidth}
        \centering
        \includegraphics[width=\linewidth]{Canada_buip_scatter.pdf}
    \end{subfigure}\hfill
    \begin{subfigure}{0.24\textwidth}
        \centering
        \includegraphics[width=\linewidth]{Euro_Area_buip_scatter.pdf}
    \end{subfigure}

    \medskip

    % ===== Row 2 =====
    \begin{subfigure}{0.24\textwidth}
        \centering
        \includegraphics[width=\linewidth]{Indonesia_buip_scatter.pdf}
    \end{subfigure}\hfill
    \begin{subfigure}{0.24\textwidth}
        \centering
        \includegraphics[width=\linewidth]{Japan_buip_scatter.pdf}
    \end{subfigure}\hfill
    \begin{subfigure}{0.24\textwidth}
        \centering
        \includegraphics[width=\linewidth]{Korea_buip_scatter.pdf}
    \end{subfigure}\hfill
    \begin{subfigure}{0.24\textwidth}
        \centering
        \includegraphics[width=\linewidth]{Mexico_buip_scatter.pdf}
    \end{subfigure}

    \medskip

    % ===== Row 3 =====
    \begin{subfigure}{0.24\textwidth}
        \centering
        \includegraphics[width=\linewidth]{New_Zealand_buip_scatter.pdf}
    \end{subfigure}\hfill
    \begin{subfigure}{0.24\textwidth}
        \centering
        \includegraphics[width=\linewidth]{Philippines_buip_scatter.pdf}
    \end{subfigure}\hfill
    \begin{subfigure}{0.24\textwidth}
        \centering
        \includegraphics[width=\linewidth]{Switzerland_buip_scatter.pdf}
    \end{subfigure}\hfill
    \begin{subfigure}{0.24\textwidth}
        \centering
        \includegraphics[width=\linewidth]{Thailand_buip_scatter.pdf}
    \end{subfigure}

    \medskip

    % ===== Row 4 =====
    \begin{subfigure}{0.24\textwidth}
        \centering
        \includegraphics[width=\linewidth]{Türkiye_buip_scatter.pdf}
    \end{subfigure}\hfill
    \begin{subfigure}{0.24\textwidth}
        \centering
        \includegraphics[width=\linewidth]{United_Kingdom_buip_scatter.pdf}
    \end{subfigure}\hfill
    \begin{subfigure}{0.24\textwidth}
        \centering
        % Empty
    \end{subfigure}\hfill
    \begin{subfigure}{0.24\textwidth}
        \centering
        % Empty
    \end{subfigure}

    \caption{BUIP model: scatter plots of fundamentalist market share $\omega_t$ against utility differential $U_{f,t} - U_{c,t}$.}
    \label{fig:buip_scatter_4x4}
\end{figure}

\subsubsection{Model Comparison}
\label{sec:model_comparison}

Table~\ref{tab:comprehensive_model_comparison} summarizes the in-sample performance of the two models under study (STR and BUIP) against a linear benchmark for each country, defined as follows:
\[
    y_t = \phi_0 + \phi_c \Delta s_{t-1} + \phi_f \eta_{t-1} + \varepsilon_t,
\]
with $\varepsilon_t \sim iid(0, \sigma^2)$.  The STR model estimation is invalid for Australia, Mexico, Thailand, Türkiye, and the UK because estimated $c$ values fall outside the data range. For the BUIP model, Japan, Korea, and the UK for the same reason. This means that the results for these countries should be interpreted with caution, or ignored. To compare the statistical performance of the three models, ill-identification cases are indicated in parentheses in Table~\ref{tab:comprehensive_model_comparison} and won't be considered.


\begin{table}[htbp]
    \centering
    \caption{Comprehensive Model Comparison: RMSE and AIC by Country}
    \label{tab:comprehensive_model_comparison}
    \begin{tabular}{l|ccc|ccc}
        \toprule
                       & \multicolumn{3}{c|}{RMSE (lower is better)} & \multicolumn{3}{c}{AIC (lower is better)}                                                                               \\
        Country        & Linear                                      & STR                                       & BUIP            & Linear            & STR               & BUIP              \\
        \midrule
        Australia      & 0.0322                                      & (0.0322)                                  & \textbf{0.0319} & \textbf{-3506.13} & (-3485.40)        & -3494.91          \\
        Brazil         & 0.1187                                      & 0.1145                                    & \textbf{0.1060} & -2180.25          & -2206.95          & \textbf{-2281.42} \\
        Canada         & 0.0205                                      & 0.0204                                    & \textbf{0.0202} & \textbf{-3981.07} & -3971.79          & -3975.31          \\
        Euro Area      & 0.0268                                      & 0.0267                                    & \textbf{0.0261} & \textbf{-2289.66} & -2270.48          & -2283.96          \\
        Indonesia      & 0.0609                                      & 0.0571                                    & \textbf{0.0442} & -2608.09          & -2644.93          & \textbf{-2888.09} \\
        Japan          & 0.0302                                      & \textbf{0.0296}                           & (0.0297)        & -3339.64          & \textbf{-3347.17} & (-3335.19)        \\
        Korea          & 0.0290                                      & \textbf{0.0275}                           & (0.0274)        & -3618.03          & \textbf{-3646.03} & (-3657.41)        \\
        Mexico         & 0.0535                                      & (0.0433)                                  & \textbf{0.0522} & -2999.03          & (-3189.84)        & \textbf{-3004.90} \\
        New Zealand    & 0.0342                                      & \textbf{0.0337}                           & 0.0337          & \textbf{-3247.73} & -3236.92          & -3242.61          \\
        Philippines    & 0.0289                                      & \textbf{0.0283}                           & 0.0287          & -3332.12          & \textbf{-3332.67} & -3319.56          \\
        Switzerland    & 0.0307                                      & 0.0304                                    & \textbf{0.0303} & \textbf{-3574.48} & -3573.40          & -3567.17          \\
        Thailand       & 0.0273                                      & (0.0274)                                  & \textbf{0.0270} & \textbf{-3113.05} & (-3075.45)        & -3102.66          \\
        Türkiye        & 0.0560                                      & (0.0529)                                  & \textbf{0.0537} & -2698.32          & (-2727.91)        & \textbf{-2720.02} \\
        United Kingdom & \textbf{0.0278 }                            & (0.0277)                                  & (0.0272)        & \textbf{-3677.91} & (-3660.45)        & (-3677.27)        \\
        \bottomrule
    \end{tabular}
    \begin{tablenotes}[para]
        \small
        \textit{Note}: Bold values indicate the best (lowest) measure for each metric within each country. Values in parentheses indicate models where the threshold/cost parameter ($c,k$) lies outside the empirical support of the transition variable, suggesting potential non-identification. For these cases, results should be interpreted with caution.
    \end{tablenotes}
\end{table}

According to RMSE, a nonlinear specification is preferred in 13 out of 14 countries: BUIP attains the lowest RMSE in nine economies (including Australia, Brazil, Canada, Euro Area, Indonesia, Switzerland, Thailand), while STR attains the lowest RMSE in three economies (Japan, Korea, New Zealand, Philippines). This confirms that, in terms of mean squared error, explicitly modeling heterogeneous trading dynamics improves in-sample fit for most of the currency pairs, with a notable improvement with the fully endogenous Behavioral UIP model. When penalizing for complexity via the AIC information criteria, linearity is rewarded in Advanced Economies. The linear benchmark is preferred in seven countries (Australia, Canada, Euro Area, New Zealand, Switzerland, Thailand, United Kingdom), illustrating that added complexity is not always rewarded once parsimony is accounted for. Among the nonlinear models, BUIP is preferred by AIC in Brazil, Indonesia, Mexico and Türkiye whereas STR is preferred in Japan, Korea and Philippines. This result suggests that nonlinear models (whether reduced form STR or the structural BUIP) matter more for emergent economies. Overall, these results suggest that nonlinearities matter for fit, but the degree to which they justify additional complexity is heterogeneous across markets and depends on their structure.

\subsubsection{Advanced vs Emerging Economies}

To synthesize the cross-country results, it is informative to contrast advanced and emerging economies along three dimensions: (i) evidence on fundamentalist forces, (ii) heterogeneity of technical trading, and (iii) the extent to which complexity is justified by the data.

First, stabilizing fundamentalist forces are identified more frequently in the BUIP framework than in the STR framework, but they are not universal in either model. In the BUIP estimates, a negative and significant $\beta_f$ emerges in several advanced economies (Australia, Euro Area, New Zealand, Switzerland) and in multiple emerging economies (Indonesia, Korea, Philippines, Thailand), suggesting that fundamental mean reversion is not confined to a particular group. Also, there is no clear distinction in terms of magnitude, so we cannot identify any particular regularity. At the same time, Mexico and Türkiye exhibit a positive and significant $\beta_f$ in BUIP, and several countries display imprecisely estimated $\beta_f$, highlighting that the empirical relevance of the PPP fundamental might be market dependent, but not systematically so across advanced vs emerging economies.

Second, chartist behavior is heterogeneous, particularly among emerging markets. BUIP estimates show significant positive chartism in Brazil, Mexico, Philippines, and Türkiye, while other countries (Canada, Indonesia, New Zealand) display significant contrarian expectations. This heterogeneity is compatible with the idea that the composition of short-horizon strategies varies across markets and periods, and that ``technical trading'' need not imply trend following.

Third, parsimony is rewarded more often in advanced economies. The AIC favors the linear benchmark in all advanced and highly liquid markets, except Japan, while nonlinear models are more often preferred using AIC in emerging markets. In contrast, RMSE results indicate that nonlinear models frequently improve fit even when AIC does not select them, suggesting that nonlinearities can matter for fit but may be rare or insufficiently persistent to justify additional parameters.

Overall, the evidence supports a view in which exchange rates can reflect nonlinear interactions between heterogeneous expectations, but the empirical manifestation of these mechanisms---and the extent to which they warrant structural complexity---differs across currency pairs.



\section{Concluding Remarks}
\label{chap:conclusions}

This paper examined whether a simple heterogeneous-expectations mechanism, with switching between fundamentalist and chartist forecasting rules, can account for monthly exchange-rate dynamics for fourteen currencies vis-\`a-vis the U.S.\ dollar over long samples ending in 2025. We estimated (i) a reduced-form logistic smooth-transition regression (STR), where regime changes are driven by observable transition variables, and (ii) a structural behavioral UIP (BUIP) specification, where regime weights are endogenized through a profitability-based discrete-choice mechanism.

Four conclusions follow from the empirical analysis. First, there is broad evidence of state dependence in exchange-rate returns. In the STR framework, linearity is rejected for 13 out of 14 currencies. Nonlinear models reduce RMSE relative to a linear benchmark in 13 out of 14 countries. At the same time, when penalizing complexity via the AIC, the linear benchmark is preferred in advanced economies, indicating that nonlinearities can be present but not always persistent or strong enough to justify additional parameters. To study emergent economies FX markets nonlinearities appear to be more relevant.

Second, the reduced-form STR provides only selective identification of the economic forces suggested by heterogeneous-agent interpretations. Evidence for trend-following behavior is limited, and PPP-consistent mean reversion is detected only for a subset of currencies; in some cases the estimated fundamentalist coefficient has a sign inconsistent with mean reversion to the PPP fundamental. These results suggest that a parsimonious STR may capture state dependence but can struggle to map it into the hypothesized behavioral channels.

Third, the structural BUIP specification of \citet{Proano2011,proano2013monetary} strengthens economic interpretation and more frequently identifies stabilizing PPP-consistent mean reversion. The fundamentalist adjustment parameter is negative and statistically significant in eight currencies (including both advanced and emerging markets), but it is not universally identified and turns significantly positive for Mexico and T\"urkiye. This pattern indicates that the empirical relevance of PPP as the fundamental value is market-dependent, and that the meaning of the fundamentalist term is sensitive to the specification of the fundamental forecast.

Fourth, technical trading is heterogeneous across currency pairs. The BUIP estimates show significant positive chartist dynamics for several currencies, while other cases exhibit significant negative coefficients, consistent with contrarian short-horizon behavior rather than trend extrapolation. This heterogeneity is consistent with the view that technical trading in FX markets encompasses multiple strategies whose dominance varies across markets and periods. Moreover, identification of the switching mechanism remains uneven: switching-intensity and threshold/cost parameters are precisely estimated only in a subset of cases, and in some currencies the estimated switching point ($c$ or $k$) lies outside the empirical support of the corresponding driving variable, calling for caution in interpretation.

Taken together, the evidence supports a view in which exchange-rate returns reflect nonlinear interactions between heterogeneous expectations, and in which such heterogeneity can improve in-sample fit. However, the gains from additional complexity are not uniform across markets, and the results highlight the need for more complex specifications mainly for emerging economies FX markets.

Several extensions appear promising. On the econometric side, allowing for time-varying volatility (e.g., STR--GARCH as in \citet{WesterhoffReitz2003,reitz2007commodity} or related heteroskedastic specifications) and exploring rolling or time-varying parameter estimation could clarify how crisis episodes affect nonlinear dynamics and switching. On the structural side, enriching the set of forecasting rules, introducing alternative fundamentals beyond absolute PPP, and allowing utilities to incorporate risk considerations or longer memory may improve identification and provide a tighter link between exchange-rate dynamics, policy regimes, and macro fundamentals.

\section*{Funding Declaration}

No funding is to be declared.



\clearpage

% ---------------------------------------------------------------------------
% Bibliography — Alphabetical order (Natbib)
% ---------------------------------------------------------------------------
%\phantomsection
\addcontentsline{toc}{section}{Bibliography}
\bibliographystyle{apalike}
\bibliography{Behavioral_FX_Dynamics.bib}

\clearpage
% ---------------------------------------------------------------------------
% Appendix
% ---------------------------------------------------------------------------

\appendix
\section{Appendix}


\subsection{Data Sources and Transformations}
\label{app:data_details}

Table~\ref{tab:data_availability} summarizes the data availability and sources for each country in the sample. The exchange rate series correspond to nominal bilateral rates vis-\`a-vis the U.S.\$, expressed as the domestic price of U.S.\$. Consumer price indices (CPI) are used to construct real exchange rates. Policy or money market rates serve as short-term interest rate proxies. All series are monthly frequency except for Canada and Australia CPI, for which we use quarterly data and interpolate to monthly frequency due to the data not being available. The data are fetched via the DBnomics API directly into python using the \texttt{dbnomics} package.

\begin{table}[htbp]
    \centering
    \caption{Data Availability and Sources by Country. S=Exchange Rate, CPI=Consumer Price Index, i=Policy/Interest Rate, MM=Money Market Rate. All data accessed via DBnomics API.}
    \label{tab:data_availability}
    \begin{tabular}{lllrl}
        \toprule
        Country        & Start Date & End Date & Observations & Source                            \\
        \midrule
        Australia      & 1982-09    & 2025-04  & 512          & IMF/IFS (S, CPI interp., i)       \\
        Brazil         & 1982-09    & 2025-06  & 514          & IMF/IFS (S, CPI, MM rate)         \\
        Canada         & 1982-09    & 2025-06  & 514          & IMF/IFS (S, CPI), OECD (i)        \\
        Euro Area      & 1999-02    & 2025-07  & 318          & IMF/IFS (S, CPI), OECD (i)        \\
        Indonesia      & 1986-06    & 2025-05  & 468          & IMF/IFS (S, CPI, i combined)      \\
        Japan          & 1985-08    & 2025-06  & 479          & IMF/IFS (S, CPI), OECD (i)        \\
        Korea          & 1982-09    & 2025-05  & 513          & IMF/IFS (S, CPI, MM rate)         \\
        Mexico         & 1982-09    & 2025-06  & 514          & IMF/IFS (S, CPI, MM rate)         \\
        New Zealand    & 1985-02    & 2025-04  & 483          & IMF/IFS (S, CPI interp., MM rate) \\
        Philippines    & 1982-09    & 2021-12  & 472          & IMF/IFS (S, CPI, MM rate)         \\
        Switzerland    & 1982-09    & 2025-07  & 515          & IMF/IFS (S, CPI), OECD (i)        \\
        Thailand       & 1989-02    & 2025-03  & 434          & IMF/IFS (S, CPI, MM rate)         \\
        Türkiye        & 1986-05    & 2025-06  & 470          & IMF/IFS (S, CPI, i)               \\
        United Kingdom & 1982-09    & 2025-07  & 515          & IMF/IFS (S, CPI), OECD (i)        \\
        \bottomrule
    \end{tabular}
\end{table}


\subsection{Model Validation Summary}
\label{app:validation_summary}

\begin{table}[htbp]

    \caption{Identification Summary: STR Threshold and BUIP Cost Parameter (In-Range Checks)}
    \label{tab:validation_summary}
    \centering
    \begin{tabular}{lcc}
        \toprule
        Country        & STR: $c$ in range of $z_t$ & BUIP: $k$ in range of $(U_f-U_c)$ \\
        \midrule
        Australia      & NO                         & YES                               \\
        Brazil         & YES                        & YES                               \\
        Canada         & YES                        & YES                               \\
        Euro Area      & YES                        & YES                               \\
        Indonesia      & YES                        & YES                               \\
        Japan          & YES                        & NO                                \\
        Korea          & YES                        & NO                                \\
        Mexico         & NO                         & YES                               \\
        New Zealand    & YES                        & YES                               \\
        Philippines    & YES                        & YES                               \\
        Switzerland    & YES                        & YES                               \\
        Thailand       & NO                         & YES                               \\
        Türkiye        & NO                         & YES                               \\
        United Kingdom & NO                         & NO                                \\
        \bottomrule
    \end{tabular}
    \begin{tablenotes}[para,flushleft]
        \small
        \textit{Note}: ``In range'' indicates whether the estimated threshold/cost parameter lies within the empirical support of the corresponding transition variable. ``NO'' flags potential identification concerns in the switching point.
    \end{tablenotes}
\end{table}



\subsection{STR model additional figures}

\begin{figure}[H] % own page in appendix
    \centering

    % ===== Row 1 =====
    \begin{subfigure}{0.30\textwidth}
        \centering
        \includegraphics[width=\linewidth]{Australia_str_detailed.pdf}
        \caption{Australia}
    \end{subfigure}\hfill
    \begin{subfigure}{0.30\textwidth}
        \centering
        \includegraphics[width=\linewidth]{Brazil_str_detailed.pdf}
        \caption{Brazil}
    \end{subfigure}\hfill
    \begin{subfigure}{0.30\textwidth}
        \centering
        \includegraphics[width=\linewidth]{Canada_str_detailed.pdf}
        \caption{Canada}
    \end{subfigure}

    \medskip

    % ===== Row 2 =====
    \begin{subfigure}{0.30\textwidth}
        \centering
        \includegraphics[width=\linewidth]{Euro_Area_str_detailed.pdf}
        \caption{Euro Area}
    \end{subfigure}\hfill
    \begin{subfigure}{0.30\textwidth}
        \centering
        \includegraphics[width=\linewidth]{Indonesia_str_detailed.pdf}
        \caption{Indonesia}
    \end{subfigure}\hfill
    \begin{subfigure}{0.30\textwidth}
        \centering
        \includegraphics[width=\linewidth]{Japan_str_detailed.pdf}
        \caption{Japan}
    \end{subfigure}

    \medskip

    % ===== Row 3 =====
    \begin{subfigure}{0.30\textwidth}
        \centering
        \includegraphics[width=\linewidth]{Korea_str_detailed.pdf}
        \caption{Korea}
    \end{subfigure}\hfill
    \begin{subfigure}{0.30\textwidth}
        \centering
        \includegraphics[width=\linewidth]{Mexico_str_detailed.pdf}
        \caption{Mexico}
    \end{subfigure}\hfill
    \begin{subfigure}{0.30\textwidth}
        \centering
        \includegraphics[width=\linewidth]{New_Zealand_str_detailed.pdf}
        \caption{New Zealand}
    \end{subfigure}

    \medskip

    % ===== Row 4 =====
    \begin{subfigure}{0.30\textwidth}
        \centering
        \includegraphics[width=\linewidth]{Philippines_str_detailed.pdf}
        \caption{Philippines}
    \end{subfigure}\hfill
    \begin{subfigure}{0.30\textwidth}
        \centering
        \includegraphics[width=\linewidth]{Switzerland_str_detailed.pdf}
        \caption{Switzerland}
    \end{subfigure}\hfill
    \begin{subfigure}{0.30\textwidth}
        \centering
        \includegraphics[width=\linewidth]{Thailand_str_detailed.pdf}
        \caption{Thailand}
    \end{subfigure}

    \medskip

    % ===== Row 5 =====
    \begin{subfigure}{0.30\textwidth}
        \centering
        \includegraphics[width=\linewidth]{Türkiye_str_detailed.pdf}
        \caption{Türkiye}
    \end{subfigure}\hfill
    \begin{subfigure}{0.30\textwidth}
        \centering
        \includegraphics[width=\linewidth]{United_Kingdom_str_detailed.pdf}
        \caption{United Kingdom}
    \end{subfigure}\hfill
    \begin{subfigure}{0.30\textwidth}
        \centering
        % Empty space
    \end{subfigure}

    \caption{STR model: Log nominal exchange rate returns (top panel), transition variable (middle panel), transition function $G_t$ (bottom panel). The dashed line in the middle panel indicates the estimated threshold $c$.}
    \label{fig:lstr_G_vs_returns_3x2}
\end{figure}

\subsection{BUIP model additional figures}

\begin{figure}[H] % own page in appendix
    \centering

    % ===== Row 1 =====
    \begin{subfigure}{0.30\textwidth}
        \centering
        \includegraphics[width=\linewidth]{Australia_buip_detailed.pdf}
        \caption{Australia}
    \end{subfigure}\hfill
    \begin{subfigure}{0.30\textwidth}
        \centering
        \includegraphics[width=\linewidth]{Brazil_buip_detailed.pdf}
        \caption{Brazil}
    \end{subfigure}\hfill
    \begin{subfigure}{0.30\textwidth}
        \centering
        \includegraphics[width=\linewidth]{Canada_buip_detailed.pdf}
        \caption{Canada}
    \end{subfigure}

    \medskip

    % ===== Row 2 =====
    \begin{subfigure}{0.30\textwidth}
        \centering
        \includegraphics[width=\linewidth]{Euro_Area_buip_detailed.pdf}
        \caption{Euro Area}
    \end{subfigure}\hfill
    \begin{subfigure}{0.30\textwidth}
        \centering
        \includegraphics[width=\linewidth]{Indonesia_buip_detailed.pdf}
        \caption{Indonesia}
    \end{subfigure}\hfill
    \begin{subfigure}{0.30\textwidth}
        \centering
        \includegraphics[width=\linewidth]{Japan_buip_detailed.pdf}
        \caption{Japan}
    \end{subfigure}

    \medskip

    % ===== Row 3 =====
    \begin{subfigure}{0.30\textwidth}
        \centering
        \includegraphics[width=\linewidth]{Korea_buip_detailed.pdf}
        \caption{Korea}
    \end{subfigure}\hfill
    \begin{subfigure}{0.30\textwidth}
        \centering
        \includegraphics[width=\linewidth]{Mexico_buip_detailed.pdf}
        \caption{Mexico}
    \end{subfigure}\hfill
    \begin{subfigure}{0.30\textwidth}
        \centering
        \includegraphics[width=\linewidth]{New_Zealand_buip_detailed.pdf}
        \caption{New Zealand}
    \end{subfigure}

    \medskip

    % ===== Row 4 =====
    \begin{subfigure}{0.30\textwidth}
        \centering
        \includegraphics[width=\linewidth]{Philippines_buip_detailed.pdf}
        \caption{Philippines}
    \end{subfigure}\hfill
    \begin{subfigure}{0.30\textwidth}
        \centering
        \includegraphics[width=\linewidth]{Switzerland_buip_detailed.pdf}
        \caption{Switzerland}
    \end{subfigure}\hfill
    \begin{subfigure}{0.30\textwidth}
        \centering
        \includegraphics[width=\linewidth]{Thailand_buip_detailed.pdf}
        \caption{Thailand}
    \end{subfigure}

    \medskip

    % ===== Row 5 =====
    \begin{subfigure}{0.30\textwidth}
        \centering
        \includegraphics[width=\linewidth]{Türkiye_buip_detailed.pdf}
        \caption{Türkiye}
    \end{subfigure}\hfill
    \begin{subfigure}{0.30\textwidth}
        \centering
        \includegraphics[width=\linewidth]{United_Kingdom_buip_detailed.pdf}
        \caption{United Kingdom}
    \end{subfigure}\hfill
    \begin{subfigure}{0.30\textwidth}
        \centering
        % Empty space
    \end{subfigure}

    \caption{BUIP model: Log nominal exchange rate returns (top panel), utility differential  (middle panel), market share of fundamentalists $\omega_t$ (bottom panel). The dashed line in the middle panel indicates the estimated cost threshold $k$.}
    \label{fig:buip_omega_vs_returns_3x2}
\end{figure}



\end{document}
